\documentclass[12pt, oneside]{article}

\usepackage{xparse}

\usepackage{geometry}
\geometry{letterpaper}
\usepackage[parfill]{parskip}
\usepackage{float}

% Math packages
\usepackage{amssymb}
\usepackage{amsmath}
\usepackage{amsthm}
\usepackage{mathalpha}

% Tables
\usepackage{booktabs}
\usepackage{afterpage}

% Hyperlinks
\usepackage{hyperref}

% References
\numberwithin{equation}{section}
\usepackage[numbers]{natbib}

% TikZ (optional for later)
\usepackage{tikz}
\usetikzlibrary{shapes,arrows,chains}

% Theorem environments
\theoremstyle{definition}
\newtheorem{definition}{Definition}[section]
\newtheorem{proposition}{Proposition}[section]
\newtheorem{axiom}{Axiom}[section]

\theoremstyle{plain}
\newtheorem{theorem}{Theorem}[section]
\newtheorem{lemma}[theorem]{Lemma}
\newtheorem{corollary}[theorem]{Corollary}
\newtheorem{example}[theorem]{Example}

\theoremstyle{remark}
\newtheorem*{remark}{Remark}

\title{Dyadic Brackets and Density-One Descent in the Collatz Map}
\author{Wayne Brassem}

\begin{document}

% Custom commands
\NewDocumentCommand{\setN}{}{\mathbb{N}}
\NewDocumentCommand{\setZ}{}{\mathbb{Z}}
\NewDocumentCommand{\setQ}{}{\mathbb{Q}}

% Document commands which provide hyperlinks to OEIS sequences referenced herein
\NewDocumentCommand{\OEIS}{m}{\href{https://oeis.org/#1}{#1}}

\maketitle

\begin{abstract}

We study the Collatz map via the fully accelerated transformation
\[
f(x)=\frac{3x+1}{2^{v_2(3x+1)}}.
\]

Finite trajectories are encoded by division words, which record the
2-adic valuations encountered along an orbit. Each such word induces
an affine map whose admissibility imposes a congruence constraint on
initial values. This yields a symbolic–affine framework in which
admissible words correspond to residue classes modulo powers of two.

We show that these constraints refine exponentially: a division word
of length $m$ restricts initial values to a residue class modulo a
power of $2$ that is comparable to, but strictly larger than, $3^m$,
while the number of admissible words grows like $\rho^m$ with $\rho < 3$.

This creates a competition between combinatorial growth and arithmetic
refinement: while the number of admissible words grows like $\rho^m$
with $\rho < 3$, each word imposes a congruence constraint modulo a
power of $2$ that grows strictly faster than $3^m$. This imbalance
forces exponential sparsification of admissible initial conditions.

As a consequence, the set of integers that fail to admit a finite
descent prefix has natural density zero. In particular, almost all
integers eventually decrease under iteration of the Collatz map.

\end{abstract}

\newpage
\tableofcontents

% ============================================================
\newpage
\section{Introduction}

The $3x+1$ problem, also known as the Collatz conjecture~\cite{Collatz1937}, concerns
the behavior of the map
\[
T(x) =
\begin{cases}
\displaystyle
    \phantom{3x} \frac{x}{2} & \text {if } x \text{ is even}, \\[8pt]
\displaystyle
    3x+1,                    & \text{if $x$ is odd}.
\end{cases}
\]
on the positive integers. It asserts that every positive integer eventually reaches $1$
under iteration of $T$. Despite its elementary formulation, the conjecture has resisted
proof for decades~\cite{Lagarias2010,Terras1976}. Extensive computational verification
confirms convergence for very large ranges~\cite{OliveiraESilva2010}, but does not
reveal the global structure governing trajectory behavior.

A major advance by Tao~\cite{Tao2019} showed that almost all integers, in the sense of
logarithmic density, eventually decrease under iteration of the Collatz map. This result
is obtained via probabilistic averaging rather than an explicit structural decomposition
of the integers.

The present work develops a deterministic, arithmetic framework for analyzing typical
Collatz behavior based on a symbolic–affine encoding of trajectories. Using this approach,
we encode the dynamics symbolically rather than studying individual trajectories with
the fully accelerated map
\[
f(x)=\frac{3x+1}{2^{v_2(3x+1)}},
\]
which maps odd integers to odd integers by compressing each odd step into a single
transformation.

Each trajectory under $f$ is represented by a finite \emph{division word}, recording the
number of factors of $2$ removed (the 2-adic valuation) at each iterate.

These words admit an exact affine representation obtained by symbolic composition of the
accelerated map: each division word induces a map of the form
\[
F_\omega(x) = \frac{3^m x + c(\omega)}{2^{D(\omega)}},
\]
and determines a corresponding congruence class modulo $2^{D(\omega)}$. Thus division words
determine canonical arithmetic \emph{constraint classes} (residue classes) in the integers,
once realizability is imposed.

This correspondence allows the dynamics to be studied at the level of residue classes
rather than individual orbits. Convergent division words—those for which the associated
affine map is contractive—identify families of integers that decrease after a fixed
number of steps. The global problem is thereby reduced to understanding how these
families are distributed and how their densities accumulate.

Two competing effects govern this distribution. On one hand, the number of admissible
convergent division words grows exponentially with length. On the other hand, each word
imposes a congruence constraint whose modulus grows exponentially with the total 2-adic
valuation. Using structural constraints on admissible division words, encoded by an
increment grammar, we show that the combinatorial growth rate is strictly subcritical:
the number of admissible words grows like $\rho^m$ for some $\rho < 3$, while the moduli
scale like powers of $2$ satisfying $3^m < 2^{D(\omega)} < 2\cdot 3^m$.

This transforms the Collatz problem into a competition between symbolic growth and
arithmetic refinement, with convergence governed by their relative exponential rates.

This imbalance leads to \emph{exponential sparsification}: the proportion of integers
admitting a convergent division word of length $m$ decays exponentially in $m$.
Summing over all lengths then shows that the complement—integers that avoid all such
descent patterns—has natural density zero.

Consequently, the set of integers whose trajectory under iteration of $f$ eventually
decreases has natural density one.

While this does not exclude the existence of exceptional trajectories, any such
exceptions must lie within a set of zero density and cannot arise from any scalable
arithmetic structure generated by the symbolic dynamics.

% \paragraph{Structure of the Paper.}
% The argument proceeds as follows:
% \begin{itemize}
% \item Section~\ref{sec:division-counts} introduces division words and their total division counts.
% \item Section~\ref{sec:affine-representation} develops the affine representation of division words and their congruence structure.
% \item Section~\ref{sec:affine-grammar} establishes admissibility and minimal valuation constraints.
% \item Section~\ref{sec:symbolic-growth} controls the growth of admissible words via the increment grammar.
% \item Section~\ref{sec:affine-constraint-fibers} identifies the residue class structure induced by division words.
% \item Section~\ref{sec:exponential-sparsification} proves exponential sparsification of admissible sets.
% \item Section~\ref{sec:density-one-descent} establishes density-one descent.
% \end{itemize}

% ============================================================
\section{Symbolic Encoding of Collatz Trajectories}
\label{sec:division-counts}
% ============================================================

% ------------------------------------------------------------
\subsection{Division Words}

\begin{definition}[Division Word]
A division word of length $m \ge 1$ is a sequence
\[
\omega = (d_0, \dots, d_{m-1}),
\]
where each $d_i \in \mathbb{N}$ denotes the number of factors of \(2\)
removed at the $i$-th iterate of the accelerated map
\[
f(x) = \frac{3x+1}{2^{v_2(3x+1)}}.
\]
\end{definition}

\begin{remark}
Since $f$ maps odd integers to odd integers, each $3f^i(x)+1$ is even,
so $d_i \ge 1$ for all $i$.
\end{remark}

\begin{definition}[Realization]
An odd integer $x \in \mathbb{N}$ realizes $\omega$ if
\[
d_i = v_2(3f^i(x)+1), \quad 0 \le i < m.
\]
\end{definition}

\begin{remark}
Every odd integer \(x\) determines a unique infinite valuation sequence
\[
(v_2(3x+1),\, v_2(3f(x)+1),\, v_2(3f^2(x)+1),\dots),
\]
of which a division word is a finite prefix.
\end{remark}

% ------------------------------------------------------------
\subsection{Total Division Count}

\begin{definition}
The total division count of a word $\omega$ is
\[
D(\omega) := \sum_{i=0}^{m-1} d_i.
\]
\end{definition}

\begin{remark}
The quantity \(D(\omega)\) measures the cumulative dyadic contraction
along the symbolic trajectory encoded by \(\omega\).
\end{remark}

% ============================================================
\section{Affine Representation of Division Words}
\label{sec:affine-representation}
% ============================================================

\begin{definition}
Each division word \(\omega\) of length \(m\) determines a formal affine map
\[
F_\omega(x)=\frac{3^m x+c(\omega)}{2^{D(\omega)}},
\]
obtained by symbolic composition of the accelerated Collatz map.
\end{definition}

\begin{proposition}
The map $\omega \mapsto F_\omega$ defines a semigroup homomorphism from
concatenation of division words (read left to right) to composition of affine maps.
\end{proposition}

\begin{proof}
Concatenating division words corresponds to successive composition
of the associated accelerated affine transformations. The exponents
add under concatenation, while the affine constants compose according
to the standard affine composition law.
\end{proof}

% ------------------------------------------------------------
\subsection{Dyadic-Ternary Balance}

The symbolic dynamics of a division word are governed by the balance
between the cumulative ternary growth \(3^m\) and the cumulative dyadic
contraction \(2^{D(\omega)}\).

For a division word \(\omega\) of length \(m\), the associated affine map
has linear coefficient
\[
\frac{3^m}{2^{D(\omega)}}.
\]

Thus:
\begin{itemize}
\item if \(2^{D(\omega)} < 3^m\), the symbolic trajectory is globally expansive,
\item if \(2^{D(\omega)} > 3^m\), the symbolic trajectory is globally contractive.
\end{itemize}

This motivates the Beatty threshold sequence
\[
A020914(m)=\lfloor m\log_2 3\rfloor+1,
\]
which is the minimal integer \(k\) satisfying
\[
2^k>3^m.
\]

Consequently, \(A020914(m)\) represents the minimal total division count
for which the affine linear coefficient becomes strictly contractive.

% ------------------------------------------------------------
\subsection{Example}

Consider convergent segments with $m=4$ iterates of $f$.  
The minimal total division count for $m=4$ is
\[
A020914(4) = 7.
\]

One realizable division word attaining this minimum is
\[
\omega = (1,1,1,4),
\]
for which
\[
D(\omega) = 1+1+1+4 = 7.
\]

Composing the accelerated map according to $\omega$ yields
\[
f_\omega(x)=\frac{3^4 x + c(\omega)}{2^7}
           = \frac{81x + 65}{128},
\]
so the canonical affine constant for this word is
\[
c(\omega)=65.
\]

The integrality condition
\[
81x + 65 \equiv 0 \pmod{128}
\]
determines a unique residue class
\[
x \equiv 15 \pmod{128}.
\]
This congruence defines a constraint fiber with modulus $2^7$.
These congruence classes will be interpreted as symbolic cylinders
associated with the division word \(\omega\).

\medskip

\noindent
\textbf{Interpretation.}
The word $\omega = (1,1,1,4)$ means:

\begin{itemize}
    \item At each of the first three iterates of $f$, $3x+1$ contains exactly one factor of $2$.
    \item At the fourth iterate of $f$, $3x+1$ contains four factors of $2$.
\end{itemize}

Under the accelerated map
\[
f(x)=\frac{3x+1}{2^{v_2(3x+1)}},
\]

Thus $d_3=4$ means that the fourth iterate divides by $2^4$ in that step.
The constant $c(\omega)$ is fixed entirely by this symbolic structure and
ensures that precisely the integers in the residue class $x \equiv 15 \pmod{128}$
produce affine form integer iterates.

This example illustrates that the behavior of trajectories is governed
entirely by the symbolic data encoded in $\omega$ and its associated
affine map. The initial values that realize a given word are precisely
those satisfying the induced congruence constraint, independent of any
particular choice of representative.

The affine representation converts symbolic trajectory data into explicit
arithmetic constraints modulo powers of two, allowing realizable trajectories
to be studied through congruence geometry rather than direct iteration.

\medskip

\noindent
In general, not every division word arises from an actual trajectory;
the condition for realizability will be formulated later as an explicit
congruence constraint, defining the class of \emph{admissible} words.

% ============================================================
\section{Affine Symbolic Admissibility Grammar}
\label{sec:affine-grammar}
% ============================================================

\begin{remark}[Division Grammar as a Dynamical Encoding]
Recall from Section~\ref{sec:division-counts} that a \emph{division word}
$\omega=(d_0,\dots,d_{m-1})$ records the exact 2-adic valuation
$d_i = v_2(3 x_i + 1)$, occurring at each iterate of the fully accelerated
Collatz map.

Unlike a purely combinatorial encoding, division words retain essential
dynamical information. Each division word $\omega=(d_0,\dots,d_{m-1})$
determines the sequence of multiplicative factors of the linear coefficient
\[
r_i = \frac{3}{2^{d_i}},
\]
so that the contractive or expansive contribution of each iterate is
preserved in the linear coefficient (with the affine term handled separately
via $c(\omega)$). This structure enables quantitative control of symbolic
trajectories via affine growth and later permits the formulation of
arithmetic realizability and descent conditions, even though individual
trajectories are collapsed under the encoding.
\end{remark}

% ------------------------------------------------------------
\subsection{Explicit Form of the Affine Constant}
\begin{lemma}[Explicit Form of the Affine Constant]
Let $\omega = (d_0,\dots,d_{m-1})$ be a division word. Then the affine map
\[
F_\omega(x) = \frac{3^m x + c(\omega)}{2^{D(\omega)}}
\]
has the explicit form
\[
c(\omega)
=
\sum_{j=0}^{m-1}
3^{m-1-j}\cdot 2^{\sum_{i=0}^{j-1} d_i},
\]
where empty sums are defined as $0$.
\end{lemma}

\begin{proof}
Iterating the accelerated map symbolically gives
\[
x_{i+1}=\frac{3x_i+1}{2^{d_i}}.
\]

Repeated substitution yields
\[
x_m
=
\frac{3^m x_0}{2^{D(\omega)}}
+
\sum_{j=0}^{m-1}
\frac{
3^{m-1-j}
}{
2^{\sum_{i=j}^{m-1} d_i}
}.
\]

Multiplying by \(2^{D(\omega)}\) gives
\[
2^{D(\omega)}x_m
=
3^m x_0
+
\sum_{j=0}^{m-1}
3^{m-1-j}
2^{\sum_{i=0}^{j-1} d_i},
\]
which identifies the affine constant \(c(\omega)\).
\end{proof}

\begin{remark}
The affine constant \(c(\omega)\) records the cumulative contribution
of all additive terms generated during symbolic iteration. Earlier
terms acquire larger dyadic weights because they experience more
subsequent divisions by \(2\), while later terms acquire larger
ternary weights through repeated multiplication by \(3\).

Explicit computational examples and a direct reconstruction of the
associated affine congruence structure are provided in the appendices.
\end{remark}

% ------------------------------------------------------------
\subsection{Stepwise Realizability}

\begin{theorem}[Affine Stepwise Identity]
\label{thm:affine-stepwise-identity}
Let $\omega = (d_0,\dots,d_{m-1})$ be a division word, and define $c(\omega)$ by
\[
c(\omega)
=
\sum_{j=0}^{m-1}
3^{m-1-j}\cdot 2^{\sum_{i=0}^{j-1} d_i}.
\]

Then for any sequence $(x_0,\dots,x_m)$ satisfying
\[
x_{i+1} = \frac{3x_i+1}{2^{d_i}}, \quad 0 \le i < m,
\]
the identity
\[
2^{D(\omega)} x_m = 3^m x_0 + c(\omega)
\]
holds.
\end{theorem}

\begin{corollary}[Stepwise Integrality Criterion]
The sequence $(x_0,\dots,x_m)$ has all intermediate values integral
under the accelerated Collatz relations if and only if
\[
3^m x_0 + c(\omega) \equiv 0 \pmod{2^{D(\omega)}}.
\]
\end{corollary}

\begin{remark}[Structural Interpretation]
Each term in $c(\omega)$ records the propagated contribution of a
single additive ``+1'' created during iteration. Earlier contributions
accumulate additional powers of $3$ through forward propagation and
additional powers of $2$ through delayed subsequent divisions by powers
of two, producing a fully time-ordered encoding of the trajectory in a
single affine constant.
\end{remark}

% ------------------------------------------------------------
\subsection{Admissibility of Division Words}

\begin{definition}[Admissibility]
\label{def:admissible-division-word}
A division word $\omega$ is admissible if the congruence
\[
3^m x + c(\omega) \equiv 0 \pmod{2^{D(\omega)}}
\]
has a solution.
\end{definition}

\begin{remark}
Admissibility is an arithmetic existence condition ensuring that the
affine map $F_\omega$ attains integer values on at least one input.
\end{remark}

% ------------------------------------------------------------
\subsection{Residue Class Representation of Division Words}

\begin{theorem}[Affine Constraint Representation]
Let $\omega$ be an admissible division word of length $m$.

Then there exists a unique residue $r_\omega \pmod{2^{D(\omega)}}$ such that
\[
x \text{ satisfies } 3^m x + c(\omega) \equiv 0 \pmod{2^{D(\omega)}}
\iff
x \equiv r_\omega \pmod{2^{D(\omega)}}.
\]
\end{theorem}

\begin{proof}
From the affine representation,
\[
F_\omega(x) = \frac{3^m x + c(\omega)}{2^{D(\omega)}},
\]
integrality is equivalent to
\[
3^m x + c(\omega) \equiv 0 \pmod{2^{D(\omega)}}.
\]

Since $3^m$ is odd, it is invertible modulo $2^{D(\omega)}$, yielding
\[
x \equiv -c(\omega) \cdot (3^m)^{-1} \pmod{2^{D(\omega)}}.
\]

This determines a unique residue modulo $2^{D(\omega)}$.

Conversely, any $x$ in this class satisfies the integrality condition.
By Theorem~\ref{thm:affine-stepwise-identity}, this is equivalent to the
existence of a trajectory whose successive iterates satisfy
\[
x_{i+1} = \frac{3x_i+1}{2^{d_i}}, \quad 0 \le i < m,
\]
and hence realize the division word $\omega$.

Thus admissible division words are in one-to-one correspondence with affine
congruence constraints modulo powers of two.
\end{proof}

\begin{corollary}
An integer $x$ realizes $\omega$ if and only if
\[
x \equiv r_\omega \pmod{2^{D(\omega)}}.
\]
\end{corollary}

\begin{remark}
The set of integers realizing $\omega$ forms a full residue class modulo
$2^{D(\omega)}$, which will later be interpreted as a symbolic cylinder
in the division-word grammar.
\end{remark}

% \begin{lemma}[Prefix Nesting of Cylinders]
% \label{lem:prefix-uniqueness}
% Let \(\omega\) and \(\omega'\) be admissible division words.

% If
% \[
% C_\omega \cap C_{\omega'} \neq \varnothing,
% \]
% then one of the words is a prefix of the other.
% \end{lemma}
\begin{lemma}[Prefix Nesting of Realization Sets]
\label{lem:prefix-uniqueness}
Let $\omega,\omega'$ be admissible division words.

If the realization sets of $\omega$ and $\omega'$ intersect, then one word
is a prefix of the other.
\end{lemma}

\begin{proof}
Every integer \(x\) determines a unique infinite valuation sequence
\[
(v_2(3x+1),\, v_2(3f(x)+1),\, v_2(3f^2(x)+1),\dots).
\]

If \(x\) lies in the realization sets of both \(\omega\) and \(\omega'\),
then both words must agree with the initial segment of this valuation sequence.

Therefore the two words cannot diverge at any position before the shorter
word terminates. Hence one word is a prefix of the other.
\end{proof}

\begin{corollary}[Disjointness at Equal Length]
Distinct admissible division words of the same length define disjoint cylinders.
\end{corollary}

\begin{proof}
If two words of equal length intersected, then by
Lemma~\ref{lem:prefix-uniqueness},
one would be a prefix of the other.
Equal lengths force equality of the words.
\end{proof}

% ============================================================
\section{Cylinders Structure of the Affine Grammar}
\label{sec:cylinders-structure}
% ============================================================

The admissible division-word grammar is naturally organized by
2-adic residue classes. Rather than forming a single global tree,
the system decomposes into a \emph{forest of prefix trees}, each
rooted at a distinct residue class determined by the affine encoding.

Each tree corresponds to the evolution of division words that share
a common initial valuation constraint, and evolves independently
under the same local extension rules.

% ------------------------------------------------------------
\subsection{Cylinder of a Division Word}
\begin{definition}[Cylinder]
Let $\omega$ be an admissible division word with associated residue
class
\[
r_\omega \pmod{2^{D(\omega)}}.
\]

The cylinder associated with $\omega$ is
\[
C_\omega
:=
\{x \in \mathbb N :
x \equiv r_\omega \pmod{2^{D(\omega)}}\}.
\]
\end{definition}

\begin{remark}[Cylinders as Tree Roots]
Each cylinder $C_\omega$ serves as the root of a local prefix tree.
Within a fixed cylinder, admissible division words correspond to
successive refinements of the same 2-adic residue class. Distinct
cylinders evolve independently, producing a disjoint \emph{forest}
of symbolic trees.
\end{remark}

% ------------------------------------------------------------
\subsection{Prefix Structure of Cylinders}

\begin{lemma}[Prefix Nesting Within a Cylinder]
Let $\omega,\omega'$ be admissible division words belonging to the
same cylinder lineage.

If
\[
C_\omega \cap C_{\omega'} \neq \varnothing,
\]
then one of the words is a prefix of the other.
\end{lemma}

\begin{proof}
Every integer $x$ determines a unique infinite valuation sequence
\[
(v_2(3x+1),\, v_2(3f(x)+1),\, v_2(3f^2(x)+1),\dots).
\]

Within a fixed cylinder, admissible words are exactly finite initial
segments of this sequence. If two such segments correspond to the same
integer, they must agree up to the length of the shorter word. Hence
one is a prefix of the other.
\end{proof}

\begin{remark}[Forest Structure]
Across different cylinders, no prefix relation is enforced. Thus the
global admissible system is not a single tree, but a disjoint union of
rooted prefix trees indexed by residue classes:
\[
\text{Admissible grammar} = \bigsqcup_{\omega \text{ root}} T_\omega.
\]

This is a \emph{forest of symbolic dynamics trees}, each governed by
identical local extension rules but evolving independently.
\end{remark}

\begin{corollary}[Disjointness at Equal Depth]
Distinct admissible division words of the same length within a fixed
cylinder are disjoint. Across different cylinders, admissibility is
incomparable and no overlap occurs within a shared residue class.
\end{corollary}

% ------------------------------------------------------------
\subsection{Cylinder Contraction Threshold}

\begin{lemma}[Minimal Contraction Threshold]
Let $\omega$ be an admissible division word of length $m$. If the
associated affine map
\[
F_\omega(x)
=
\frac{3^m x+c(\omega)}{2^{D(\omega)}}
\]
is contractive in its linear coefficient, then
\[
D(\omega)\ge A020914(m),
\]
where
\[
A020914(m)=\lceil m\log_2 3\rceil
\]
is the minimal integer \(k\) satisfying
\[
2^k>3^m.
\]
\end{lemma}

\begin{proof}
Contractivity requires
\[
\frac{3^m}{2^{D(\omega)}}<1.
\]
which is equivalent to
\[
2^{D(\omega)}>3^m.
\]

The minimal integer satisfying this inequality is exactly
$A020914(m)$, so
\[
D(\omega)\ge A020914(m).
\]
\end{proof}

% ------------------------------------------------------------
\subsection{Contractive Cylinders}

\begin{definition}[Contractive Cylinder]
An admissible cylinder $C_\omega$ is contractive if
\[
\frac{3^m}{2^{D(\omega)}}<1.
\]
Equivalently,
\[
D(\omega)>m\log_2 3.
\]
\end{definition}

\begin{remark}
The Beatty threshold sequence \OEIS{A020914} defines the boundary between
expansive and contractive regions of each tree in the forest. It governs
the growth/decay balance along admissible branches within each cylinder,
with identical dynamics across the forest of admissible cylinders, independent
of the chosen root.
\end{remark}

% % ============================================================
% \section{Reconstruction of Integers from a Division Word}
% \label{sec:reconstruction}

% \subsection{Inverse Cylinder Construction}

% Let $\omega = (d_0,\dots,d_{m-1})$ be an admissible division word with total division count
% \[
% D(\omega) = \sum_{i=0}^{m-1} d_i.
% \]

% From Section~\ref{sec:affine-grammar}, $\omega$ induces the affine map
% \[
% F_\omega(x) = \frac{3^m x + c(\omega)}{2^{D(\omega)}}.
% \]

% Admissibility implies the congruence condition
% \[
% 3^m x + c(\omega) \equiv 0 \pmod{2^{D(\omega)}}.
% \]

% Since $3^m$ is invertible modulo $2^{D(\omega)}$, this yields the unique residue class
% \[
% x \equiv -c(\omega)\,(3^m)^{-1} \pmod{2^{D(\omega)}}.
% \]

% \subsection{Canonical Parametrization of the Cylinder}

% Let $r_\omega$ denote the canonical representative of this residue class in
% $\{0,1,\dots,2^{D(\omega)}-1\}$. Then all integers realizing $\omega$ are given by
% \[
% x(k) = r_\omega + k\cdot 2^{D(\omega)}, \quad k \in \mathbb{N}_0.
% \]

% \subsection{Affine Reconstruction Form}

% Applying $F_\omega$ to $x(k)$ yields
% \[
% F_\omega(x(k))
% = \frac{3^m (r_\omega + k2^{D(\omega)}) + c(\omega)}{2^{D(\omega)}}.
% \]

% Expanding,
% \[
% F_\omega(x(k))
% = k\cdot 3^m + \frac{3^m r_\omega + c(\omega)}{2^{D(\omega)}}.
% \]

% Define the constant
% \[
% K_\omega := \frac{3^m r_\omega + c(\omega)}{2^{D(\omega)}}.
% \]

% Then the affine form becomes
% \[
% F_\omega(x(k)) = k\cdot 3^m + K_\omega.
% \]

% \subsection{Interpretation}

% Thus each division word induces a bijection between:
% \begin{itemize}
% \item integers in a residue class $x(k)$, and
% \item affine images of the form $k\cdot 3^m + K_\omega$.
% \end{itemize}

% This provides a complete inverse description of the division-word cylinder.

% \subsection{Example}
% Solving the congruence
% \[
% 3^4 x + 65 \equiv 0 \pmod{128}
% \]
% gives
% \[
% 81x + 65 \equiv 0 \pmod{128}.
% \]

% Reducing $81 \equiv -47 \pmod{128}$, we obtain
% \[
% -47x + 65 \equiv 0 \pmod{128}
% \quad \Longrightarrow \quad
% 47x \equiv 65 \pmod{128}.
% \]

% Since $47$ is invertible modulo $128$, with
% \[
% 47^{-1} \equiv 79 \pmod{128},
% \]
% we multiply both sides to obtain
% \[
% x \equiv 79 \cdot 65 \equiv 15 \pmod{128}.
% \]

% Thus the division word $\omega = (1,1,1,4)$ generates the residue class
% \[
% x = 15 + 128k, \quad k \in \mathbb{N}_0.
% \]


% ============================================================
\section{Jump-Induced Admissible Grammar Expansion}
\label{sec:jump-expansion}
% ============================================================

The prefix-tree structure of admissible cylinders induces a corresponding
symbolic grammar on division words. The admissible extensions of a word
are governed by cumulative dyadic growth constraints, determined by the
minimal contraction threshold sequence \OEIS{A020914}.

The branching structure of minimally contractive division words is governed
by the prefix-sum constraints induced by
\[
B(n)=A020914(n)=\lfloor n\log_2 3\rfloor+1.
\]

This sequence plays a dual role: it is simultaneously the global Beatty
threshold for dyadic contraction and the local admissibility boundary
governing extension of division words.  For a division word
\[
\omega=(d_0,\dots,d_{m-1}),
\]
define the prefix sums
\[
H(i)=\sum_{j=0}^{i} d_j.
\]

The quantity \(H(i)\) records the cumulative dyadic exponent accumulated
along the first \(i+1\) iterates of the symbolic trajectory. The Beatty
threshold \(B(i+1)\) is the minimal exponent for which the corresponding
linear coefficient becomes contractive:
\[
2^{B(i+1)}>3^{i+1}.
\]

Thus the inequalities
\[
H(i)<B(i+1)
\quad (0\le i<m-1),
\]
ensure that every proper prefix remains strictly below the contractive
threshold, while the terminal equality
\[
H(m-1)=B(m).
\]
forces the word to achieve minimal contraction exactly at length \(m\).
These inequalities define a constrained lattice region inside
\[
\mathbb{N}^m,
\]
whose integer points correspond to minimally contractive admissible
division words of length \(m\).

\begin{definition}[Prefix-Admissible Word]
A division word $\omega=(d_0,\dots,d_{m-1})$ is admissible if and only if
its prefix sums satisfy
\[
H(i) < B(i+1) \quad (0 \le i < m-1),
\qquad
H(m-1)=B(m).
\]
\end{definition}

% ------------------------------------------------------------
\subsection{Propagation Versus Expansion}

The increment sequence
\[
\delta(n)=B(n+1)-B(n)
\]
takes values in \(\{1,2\}\).

The distinction between these two increment types governs the
combinatorial evolution of the admissible grammar.  In cylinder
language, these increments correspond to whether refinement
occurs within an existing residue class or splits it into new
congruence classes

\begin{lemma}[Unit Increments Preserve the Existing Grammar]
\label{lem:single-no-new}
Suppose
\[
B(n+1)-B(n)=1.
\]

introduces no new interior admissibility directions beyond propagation
of existing prefixes. The admissible words at level \(n+1\) arise solely
from propagation of previously admissible prefix configurations.
\end{lemma}

\begin{proof}[Structural mechanism]
A unit increment increases the terminal admissibility bound by exactly
one:
\[
B(n+1)=B(n)+1.
\]

Since every division word component satisfies \(d_i\ge1\), the final
coordinate is forced to absorb precisely this additional unit.

No new slack is introduced into the earlier prefix constraints
\[
H(i)<B(i+1),
\]
so the admissible lattice region does not acquire any new interior
directions.

Thus the grammar propagates existing admissible configurations but
creates no genuinely new combinatorial branches.
\end{proof}

\begin{lemma}[Jump Increments Create New Admissible Words]
\label{lem:jump-new-words}
Suppose
\[
B(n+1)-B(n)=2.
\]

Then extending the admissible search from length \(n\) to length
\(n+1\) necessarily creates new admissible division words that did
not occur at earlier lengths.
\end{lemma}

\begin{proof}[Geometric mechanism]
A jump increment enlarges the terminal admissibility bound by two:
\[
B(n+1)=B(n)+2.
\]

After accounting for the mandatory positivity condition
\[
d_i\ge1,
\]
one additional unit of slack remains available.

This additional slack creates new admissible lattice points satisfying
\[
H(i)<B(i+1)
\]
for all intermediate prefixes while still achieving the terminal
constraint
\[
H(n)=B(n+1).
\]

Consequently the admissible lattice region acquires new integer points,
corresponding to genuinely new admissible division words.
\end{proof}

% ------------------------------------------------------------
\subsection{First-Passage Cylinder Grammar}

\begin{theorem}[First-Passage Characterisation]
A division word is admissible if and only if its prefix-sum trajectory
achieves first passage of the Beatty boundary \(B(n)\).
\end{theorem}

\begin{corollary}[Forest Uniqueness]
Each admissible division word corresponds to a unique node in the
admissible cylinder forest, and hence to a unique cylinder.

No admissible word can occur at more than one depth.
\end{corollary}

\begin{theorem}[Forest Enumeration Principle]
The admissible cylinders form a rooted forest under prefix extension,
where each node corresponds to a unique admissible division word.

Therefore the number of admissible words of length \(m\) equals the
number of nodes at depth \(m\) across the forest.
\end{theorem}

\begin{remark}[Grammar as a Constrained Branching Forest]
The admissible grammar is therefore a rooted branching forest whose
local branching structure at depth \(n\) is determined entirely by the
increment
\[
\delta(n)=B(n+1)-B(n)\in\{1,2\}.
\]

Unit increments propagate existing admissible branches, while jump
increments create new admissible branching directions within the forest.
\end{remark}

% ============================================================
\section{Symbolic Growth and Spectral Bounds}
\label{sec:symbolic-growth}
% ============================================================

The symbolic dynamics naturally separate into two interacting
structures:

\begin{enumerate}
\item the \emph{survivor forest}, consisting of cylinders whose
prefix trajectories remain strictly below the Beatty contraction
boundary, and

\item the \emph{first-passage forest}, consisting of cylinders that
reach the contraction boundary for the first time at a specified depth.
\end{enumerate}

The survivor forest evolves by successive refinement of admissible
prefixes, while the first-passage forest records the first entry of
surviving trajectories into the contractive regime.

Throughout this section, we study the combinatorial growth of the
first-passage forest.

Let
\[
N(m)
\]
denote the number of first-passage division words of length \(m\),
equivalently the number of contractive cylinders created at depth
\(m\). By the first-passage characterization of admissible words,
\[
N(m)=A186009(m+1).
\]

% ------------------------------------------------------------
\subsection{Survivor refinement and first-passage decomposition}

For each depth \(m\), let
\[
S_m
\]
denote the set of integers whose symbolic trajectories remain
admissible through depth \(m\), equivalently, whose prefix sums
remain strictly below the Beatty boundary at all intermediate depths.
Thus
\[
S_{m+1}\subseteq S_m.
\]

\begin{remark}
The symbolic dynamics therefore define a filtration of the odd
integers by nested survivor sets together with successive escape
layers:
\[
S_1 \supseteq S_2 \supseteq S_3 \supseteq \cdots.
\]

At each depth, surviving cylinders either:
\begin{itemize}
\item refine into deeper survivor cylinders, or
\item exit into the first-passage forest.
\end{itemize}
\end{remark}

% ------------------------------------------------------------
\subsection{Cylinder concatenation principle}

Surviving symbolic branches concatenate only when their associated
cylinders are compatible.

More precisely, let \(\omega\) and \(\eta\) be surviving division
words of lengths \(m\) and \(n\). Any surviving extension of total
length \(m+n\) must lie in
\[
C_\omega \cap f^{-m}(C_\eta),
\]
where \(f^{-m}\) denotes the backward symbolic shift induced by the
valuation dynamics.

By the prefix-nesting property of cylinders, this intersection is
either empty or a single refined cylinder determined by a unique
surviving extension of \(\omega\).

Consequently, refinement of survivor cylinders produces only finitely
many admissible continuations, with branching uniformly controlled
independently of depth.

% ------------------------------------------------------------
\subsection{Submultiplicative growth}

\begin{lemma}[Submultiplicative Growth]
There exists \(C>0\) such that
\[
N(m+n)\le C\,N(m)N(n).
\]
\end{lemma}

\begin{proof}
Each first-passage word of length \(m+n\) determines uniquely:
\begin{itemize}
\item a surviving prefix of length \(m\), and
\item a compatible continuation of length \(n\) leading to first passage.
\end{itemize}

By the cylinder structure, every surviving prefix corresponds to a
residue class modulo \(2^{D(\omega)}\), and subsequent continuations
arise only through compatible refinements of this class.

The jump/propagation structure of the Beatty boundary \(B(n)\)
permits only finitely many admissible refinements at each stage.
Hence the number of compatible continuations per surviving cylinder
is uniformly bounded independently of depth.

Therefore the total number of first-passage words of length \(m+n\)
is bounded by a constant multiple of the product of the counts at
lengths \(m\) and \(n\).
\end{proof}

% ------------------------------------------------------------
\subsection{Spectral growth rate}

\begin{theorem}[Existence of Exponential Growth Rate]
The limit superior
\[
\lambda
=
\limsup_{m\to\infty} N(m)^{1/m}
\]
is finite.
\end{theorem}

\begin{proof}
The result follows from submultiplicativity via standard
Fekete-type arguments.
\end{proof}

% ------------------------------------------------------------
\subsection{Entropy versus dyadic contraction}

The quantity \(N(m)\) measures the combinatorial complexity of the
first-passage forest at depth \(m\), while each first-passage
cylinder carries dyadic density approximately
\[
2^{-D(\omega)}.
\]

For minimally contractive cylinders,
\[
D(\omega)\approx m\log_2 3,
\]
so a typical cylinder has density approximately
\[
3^{-m}.
\]

Thus the asymptotic behaviour of the symbolic dynamics is governed
by the competition between:
\begin{itemize}
\item exponential branching growth in the number of surviving
symbolic trajectories, and
\item exponential dyadic contraction in cylinder density.
\end{itemize}

The spectral quantity \(\lambda\) measures the effective entropy
production of the survivor forest. If this symbolic growth is
dominated by dyadic contraction, then the total surviving density
necessarily collapses under successive refinement.

% ============================================================
\section{Measure-Theoretic Collapse of the Surviving Forest}
\label{sec:measure-collapse}
% ============================================================

The admissible forest constructed in the previous sections naturally
induces a nested family of measurable subsets of the odd integers.
These sets encode the integers whose symbolic trajectories remain
admissible to increasing depths.

The central question is whether the total density of surviving
admissible cylinders persists indefinitely, or collapses under repeated
refinement.

% ------------------------------------------------------------
\subsection{Surviving cylinder sets}

\begin{definition}[Depth-$m$ surviving set]
Let
\[
S_m
:=
\bigcup_{|\omega|=m} C_\omega,
\]
where the union ranges over all admissible division words of length \(m\).
\end{definition}

\begin{remark}
The set \(S_m\) consists of all integers whose symbolic trajectory
remains admissible through depth \(m\). Since admissible cylinders
refine under word extension,
\[
S_{m+1}\subseteq S_m.
\]
Thus the surviving sets form a nested decreasing sequence.
\end{remark}

\subsection{First-passage layers}

\begin{definition}[First-passage layer]
Define the first-passage layer at depth \(m\) by
\[
F_m := S_m \setminus S_{m+1}.
\]
\end{definition}

\begin{remark}
The set \(F_m\) consists of integers whose symbolic trajectories
remain admissible through depth \(m\), but fail admissibility
at the next refinement step.
\end{remark}

\begin{lemma}[Disjoint survival decomposition]
For every \(m\ge1\),
\[
S_m = F_m \sqcup S_{m+1},
\]
where \(\sqcup\) denotes disjoint union.
\end{lemma}

\begin{proof}
By definition,
\[
F_m = S_m \setminus S_{m+1}.
\]
Hence every element of \(S_m\) either:
\begin{itemize}
\item survives to depth \(m+1\), or
\item exits the survivor forest at depth \(m\).
\end{itemize}

These two possibilities are mutually exclusive, giving the disjoint decomposition.
\end{proof}

\begin{remark}[Conservation of symbolic mass]
The decomposition
\[
S_m = F_m \sqcup S_{m+1}
\]
shows that symbolic refinement conserves total mass: cylinders are not
destroyed under refinement, but redistributed between deeper survivor
cylinders and first-passage escape cylinders.
\end{remark}

\begin{corollary}[Global first-passage decomposition]
\label{cor:global-decomposition}
The initial survivor set admits the disjoint decomposition
\[
S_1
=
\left(\bigsqcup_{m\ge1} F_m\right)
\sqcup
S_\infty,
\]
where
\[
S_\infty
=
\bigcap_{m\ge1} S_m
\]
is the infinite surviving set.
\end{corollary}

\begin{proof}
Iterating the decomposition
\[
S_m = F_m \sqcup S_{m+1}
\]
gives
\[
S_1
=
F_1 \sqcup F_2 \sqcup \cdots \sqcup F_m \sqcup S_{m+1}.
\]

Passing to the limit \(m\to\infty\) yields
\[
S_1
=
\left(\bigsqcup_{m\ge1}F_m\right)
\sqcup
S_\infty.
\qedhere
\]
\end{proof}

\begin{remark}[Symbolic conservation law]
The global decomposition shows that every surviving trajectory
either:
\begin{itemize}
\item eventually exits into a first-passage layer, or
\item survives indefinitely in \(S_\infty\).
\end{itemize}

No symbolic mass is lost or duplicated under refinement; it is
partitioned disjointly between finite first-passage escape layers
and the infinite surviving set.
\end{remark}

% ------------------------------------------------------------
\subsection{Cylinder measure}

\begin{definition}[Natural dyadic measure]
Let \(C_\omega\) be a cylinder determined by a residue class modulo
\(2^{D(\omega)}\). Define its measure by
\[
\mu(C_\omega)
=
2^{-D(\omega)}.
\]
\end{definition}

\begin{remark}
Since each admissible cylinder corresponds to a single residue class
modulo \(2^{D(\omega)}\), its canonical dyadic weight is proportional
to \(2^{-D(\omega)}\). Up to normalization of the odd integers, this
coincides with the dyadic scaling above.
\end{remark}

\begin{lemma}[Disjointness at fixed depth]
Distinct admissible cylinders of equal depth are disjoint.
\end{lemma}

\begin{proof}
This follows from the prefix-nesting property established in
Section~\ref{sec:cylinders-structure}.
\end{proof}

\begin{corollary}[Measure decomposition]
For every \(m\),
\[
\mu(S_m)
=
\sum_{|\omega|=m}2^{-D(\omega)}.
\]
\end{corollary}

\begin{remark}[Conservation under symbolic refinement]
The first-passage decomposition partitions the initial surviving
set into disjoint escape layers together with the infinite
surviving set.

Thus dyadic mass is neither created nor lost under refinement:
it is redistributed between progressively finer symbolic layers.
\end{remark}

% ------------------------------------------------------------
\subsection{Nested survival structure}

\begin{definition}[Infinite surviving set]
Define the infinite surviving set by
\[
S_\infty
:=
\bigcap_{m=1}^{\infty} S_m.
\]
\end{definition}

\begin{remark}
The set \(S_\infty\) consists of integers whose symbolic trajectories
remain admissible at every depth, equivalently, integers which never
exit the admissible forest under successive refinement.
\end{remark}

\begin{theorem}[Continuity of measure under nested refinement]
\label{thm:nested-measure}
Since
\[
S_{m+1}\subseteq S_m,
\]
the measures satisfy
\[
\mu(S_\infty)
=
\lim_{m\to\infty}\mu(S_m).
\]
\end{theorem}

\begin{proof}
This is the standard continuity property of measure for nested
decreasing measurable sets.
\end{proof}

% ------------------------------------------------------------
\subsection{Entropy versus dyadic contraction}

The surviving mass at depth \(m\) is governed by two competing effects:

\begin{enumerate}
\item combinatorial growth in the number of admissible cylinders,
\item dyadic decay in the measure of each cylinder.
\end{enumerate}

If \(N(m)\) denotes the number of admissible words of length \(m\),
then formally
\[
\mu(S_m)
=
\sum_{|\omega|=m}2^{-D(\omega)}.
\]

For minimally contractive admissible words,
\[
D(\omega)\approx m\log_2 3,
\]
so individual cylinders have typical size
\[
2^{-D(\omega)}
\approx
3^{-m}.
\]

Meanwhile the number of admissible cylinders grows exponentially at rate
\[
\lambda
=
\limsup_{m\to\infty}N(m)^{1/m}.
\]

Consequently the surviving mass heuristically satisfies
\[
\mu(S_m)
\approx
\left(\frac{\lambda}{3}\right)^m.
\]

% ------------------------------------------------------------
\subsection{Spectral collapse criterion}

\begin{theorem}[Heuristic spectral collapse criterion]
\label{thm:spectral-collapse}
If
\[
\lambda < 3,
\]
then
\[
\mu(S_m)\to0.
\]
Consequently,
\[
\mu(S_\infty)=0.
\]
\end{theorem}

\begin{proof}[Proof sketch]
The exponential growth of admissible cylinders is dominated by the
dyadic decay of cylinder measure. Formally,
\[
\mu(S_m)
=
\sum_{|\omega|=m}2^{-D(\omega)}
\lesssim
N(m)\cdot 3^{-m}.
\]

If
\[
N(m)\lesssim \lambda^m
\]
with \(\lambda<3\), then
\[
\mu(S_m)\lesssim \left(\frac{\lambda}{3}\right)^m,
\]
which converges exponentially to zero.
The conclusion follows from
Theorem~\ref{thm:nested-measure}.
\end{proof}

\begin{remark}
The remaining task is therefore reduced to obtaining sufficiently sharp
spectral bounds on the admissible grammar to establish
\[
\lambda<3.
\]

This transforms the global density problem into a symbolic growth
problem on the admissible forest.
\end{remark}

% ------------------------------------------------------------
\subsection{Interpretation}

The admissible forest therefore behaves as a symbolic filtration of the
odd integers. Increasing depth imposes increasingly restrictive dyadic
constraints, refining the integers into thinner admissible cylinders.

The balance between:
\begin{itemize}
\item symbolic branching entropy, and
\item dyadic cylinder contraction
\end{itemize}
determines whether positive density can survive indefinitely.

If dyadic contraction dominates symbolic growth, then the total
surviving density collapses to zero, leaving only a null exceptional
set of eternally surviving trajectories.

% ============================================================
\section{Extremal Growth via Sturmian Return Words}
% ============================================================

\subsection{Sturmian Structure of the Dyadic Threshold Sequence}

Let
\[
\alpha := \log_2 3,
\qquad
d(n) := \lfloor n\alpha \rfloor + 1.
\]

The increment sequence
\[
\delta(n) := d(n+1) - d(n) \in \{1,2\}
\]
is a Sturmian word of slope $\alpha - 1$ and intercept $0$.

Equivalently, $\delta$ defines a symbolic dynamical system
\[
X_\alpha \subset \{1,2\}^{\mathbb{N}}
\]
generated by an irrational rotation.

All admissible growth structures considered in this paper are functions
of this Sturmian coding.

% ------------------------------------------------------------

\subsection{Return Words and Primitive Cycle Decomposition}

\begin{definition}[Return Words]
A finite word $C$ is a return word of the Sturmian system $X_\alpha$
if it is the minimal block between successive visits to a fixed cylinder
set induced by $\delta$.

We denote by $\mathcal{R}_\alpha$ the set of return words.
\end{definition}

\begin{theorem}[Return Words via Continued Fractions]
Let $\alpha = \log_2 3$ and let $\delta(n)$ be the Sturmian coding induced
by rotation $R_\alpha$.

Let $q_k$ denote the denominators of the continued-fraction convergents of $\alpha$.

Then for any cylinder set defining the admissible grammar, the set of return
words is exactly:
\[
\mathcal{R}_\alpha = \{C_{q_k}, C_{q_{k+1}}\}
\]
for some index $k$ determined by the chosen prefix.

In particular, for the admissible threshold structure of $d(n)=\lfloor n\alpha \rfloor +1$,
the relevant convergent pair is:
\[
(q_k, q_{k+1}) = (5,12).
\]
\end{theorem}

\begin{proof}
The sequence $\delta(n)$ is a Sturmian word of slope $\alpha$.
Return words to any prefix correspond to first-return maps of the rotation
$R_\alpha$ to a subinterval $I$.

By standard Sturmian return-word theory, the induced first-return map is an
exchange of two intervals, and the return times are exactly the two consecutive
denominators of continued-fraction convergents of $\alpha$.

Since $\alpha = \log_2 3$ has continued fraction expansion
\[
\alpha=[1;1,1,2,2,3,1,5,2,23,2,2,\dots].
\]
its convergent denominators begin
\[
1,\ 1,\ 2,\ 5,\ 12,\ 41,\ 53,\ 306,\ 665,\ 15601,\ldots
\]

The first convergent pair whose associated cocycle evaluations
exhibit opposite drift signs occurs at the transition
\[
5 \to 12.
\]

Hence the return-word set is exactly $\{C_5, C_{12}\}$.
\end{proof}

\begin{remark}[Diophantine Origin of the Grammar]
\label{rem:diophantine-grammar}

The admissible grammar is fundamentally Diophantine rather than
purely combinatorial.

The convergent denominators
\[
1,\ 1,\ 2,\ 5,\ 12,\ 41,\ 53,\ 306,\ 665,\ 15601,\ldots
\]
arise from the continued-fraction approximants of
\[
\alpha=\log_2 3,
\]
and therefore correspond to the best rational approximations
\[
\frac{p_k}{q_k}\approx \log_2 3.
\]

Equivalently,
\[
2^{p_k}\approx 3^{q_k},
\]
so the associated multiplicative cocycle satisfies
\[
\frac{2^{p_k}}{3^{q_k}}\to1.
\]

The observed alternating expansive and contractive tower structure is
therefore a direct manifestation of the alternating sign structure of
continued-fraction convergents.

In particular, the appearance of the 5-, 7-, 12-, 41-, and 53-cycles
is not accidental; these cycles are induced by the Diophantine
approximation hierarchy of $\log_2 3$.
\end{remark}

\begin{lemma}[Induced Return-Word Hierarchy]
\label{lem:induced-hierarchy}

Every admissible word in the Sturmian language generated by
\[
\delta(n)=\lfloor (n+1)\alpha\rfloor-\lfloor n\alpha\rfloor
\]
admits a hierarchical decomposition into return words associated with
successive convergent scales of $\alpha$.

In particular, all higher return structures are generated recursively
from the primitive pair
\[
(C_5,C_{12}),
\]
corresponding to the first nontrivial convergent denominators
\[
(5,12).
\]

Consequently, longer admissible blocks such as the
12-, 41-, and 53-cycles introduce no new primitive local growth
mechanisms beyond those already present in $C_5$ and $C_{12}$.
\end{lemma}

\noindent
This reduces all global growth questions to finite combinations of two primitive return words.

% ------------------------------------------------------------

\subsection{Multiplicative Cocycle Structure}

Define the multiplicative growth cocycle
\[
G(w) := \frac{2^{\#s(w) + 2\#d(w)}}{3^{|w|}},
\]
where $\#s(w)$ and $\#d(w)$ denote the number of single and double increments in $w$.

\begin{lemma}[Cocycle Additivity]
For concatenation of admissible return words,
\[
G(C_{i_1} \cdots C_{i_n}) = \prod_{k=1}^n G(C_{i_k}).
\]
Equivalently,
\[
\log G(w) = \sum_{k} \log G(C_{i_k}).
\]
\end{lemma}

Thus the global exponential growth rate reduces to weighted averages
of return-word contributions.

% ------------------------------------------------------------

\subsection{Local Evaluation of Return Words}

The two primitive return words satisfy:

\[
\frac{1}{5}\log G(C_5)
=
\log\!\left(
3^{2}
\cdot \frac{123}{33}
\cdot \frac{341}{131}
\cdot \frac{329}{123}
\right)^{1/5},
\]

\[
\frac{1}{7}\log G(C_7)
<
\frac{1}{5}\log G(C_5).
\]

Hence $C_5$ is the unique extremal return word with maximal per-symbol drift.

\begin{proposition}[Extremal Return Word Principle]
For any admissible infinite word $\delta \in X_\alpha$,
\[
\limsup_{n\to\infty}
\frac{1}{n} \log G(\delta_{[1,n]})
\le
\frac{1}{5}\log G(C_5).
\]
\end{proposition}

\begin{proof}
By the unique decomposition into return words and cocycle additivity,
the growth rate is a convex combination of the per-symbol contributions
of $C_5$ and $C_7$.

Since
\[
\frac{1}{5}\log G(C_5)
>
\frac{1}{7}\log G(C_7),
\]
the maximal average is achieved by the periodic word $C_5^\infty$.
\end{proof}

% ------------------------------------------------------------

\subsection{Global Growth Bound}

\begin{theorem}[Spectral Growth Bound]
Let $\lambda$ denote the exponential growth rate of admissible configurations
under the horizontal-sum dynamics.

Then
\[
\lambda
\le
\exp\!\left(\frac{1}{5}\log G(C_5)\right)
< 3.
\]
\end{theorem}

\begin{proof}
From the extremal return-word principle, the maximal growth is achieved by
periodic concatenation of $C_5$. Evaluating the local constants yields
\[
\left(
3^{2}
\cdot \frac{123}{33}
\cdot \frac{341}{131}
\cdot \frac{329}{123}
\right)^{1/5}
< 3.
\]
\end{proof}

% ------------------------------------------------------------

\subsection{Interpretation}

The admissible grammar is therefore governed by a Sturmian shift with two
primitive return words. Although both return words occur in all admissible
sequences, the 5-cycle is the unique maximizer of the cocycle drift.

All longer structures (including 12-, 41-, and 53-cycles) arise as
concatenations of these primitives and do not introduce new extremal growth.

Consequently, the asymptotic growth is controlled entirely by the
5-cycle envelope, which is strictly subcritical relative to $3$.

% ============================================================
\section{Extremal Local Growth Analysis}
% ============================================================

The following section demonstrates that the maximal local growth factors
associated with admissible patterns are uniformly bounded below the critical
threshold of 3, and that the minimal admissible expansive block (the 5-cycle)
has a maximal geometric growth rate strictly below 3, ensuring that the overall
spectral growth of the admissible grammar is insufficient to overcome dyadic
contraction.

\subsection{Explicit enumeration of admissible patterns}
\begin{lemma}[Local Growth Bounds]
\label{lem:local-growth-bounds}

For each admissible local pattern occurring within \OEIS{A022921},
the ratio of the total sum of entries after applying the pattern
to the total sum before applying it is bounded above as follows,
uniformly over all admissible placements of (delayed) doubling within
the pattern:
\[
\begin{array}{c|c}
\text{Pattern} & \text{Max Growth Factor} \\
\hline
s,d & 3 \\
d,d & 123/33 \\
s,d,s & 341/131 \\
d,d,s & 329/123
\end{array}
\]
Here $s$ denotes a single increment and $d$ a double increment in the
horizontal-sum construction.  Note that these bounds are independent
of the starting row and only depend on the local pattern.
\end{lemma}

\begin{proposition}[Extremal 5-Cycle Growth]
\label{prop:extremal-5-cycle}
The maximal geometric growth rate achievable by any admissible infinite
word under the horizontal-sum dynamics associated with \OEIS{A186009} is bounded by
\[
\rho_{5}
=
\Bigl(
3^{2}
\cdot (123/33)^1
\cdot (341/131)^1
\cdot (329/123)^1
\Bigr)^{1/5}
\approx 2.976
< 3,
\]
where the exponents record the number of times each admissible local pattern
appears within a single admissible $5$-cycle (the minimal admissible expansive
block).  See Section~\ref{sec:admissible-patterns} for explicit enumeration:
two occurrences of $\{s,d\}$, one of $\{d,d\}$, one of $\{s,d,s\}$, and one
of $\{d,d,s\}$.

% This enumeration ensures that every admissible infinite word can be decomposed
% into concatenations of $5$--cycles whose aggregate growth is controlled by
% \(\rho_5\). See Lemma~\ref{lem:inevitable-5-cycle} for a rigorous argument
% that 5-cycles cannot be avoided and therefore this bound applies globally.
% \end{proposition}

% \begin{lemma}[Inevitable 5-Cycle Insertion]
% \label{lem:inevitable-5-cycle}
% Every contractive sequence of cycles under the admissibility grammar of \OEIS{A022921}
% necessarily contains at least one embedded 5-cycle within either a 12-cycle (7-cycle + 5-cycle) 
% or as a standalone 5-cycle following sufficient consecutive 12-cycles. 

% In particular:
% \begin{enumerate}
%     \item Purely consecutive 7-cycles cannot occur.
%     \item Expansive 5-cycles appear at bounded intervals within any admissible infinite word.
%     \item After sequences of three or four consecutive 12-cycles (yielding 41- or 53-cycle blocks),
%           the grammar forces a 5-cycle insertion without exceeding the contractive bound of unity.
%     \item It is therefore impossible to indefinitely avoid 5-cycles; they are structurally unavoidable.
% \end{enumerate}
% \end{lemma}

\begin{remark}[Extremal 5-Cycle Envelope]
Although the 12-cycle is contractive overall, it contains a 5-cycle in its tail. 
Repeated sequences of 12-cycles may occasionally yield consecutive 5-cycles once 
enough contractive slack has accumulated, as seen in the 41-cycle (3 × 12 + 5) 
or 53-cycle (4 × 12 + 5) structures, but the grammar ensures that 5-cycles 
cannot be skipped indefinitely. 

This observation justifies treating an infinite repetition of 5-cycles as a strict upper envelope: 
it saturates all local growth factors, while any admissible word, including \OEIS{A186009}, 
contains interspersed contractive 12-cycles, so its per-symbol growth is strictly below 
the 5-cycle maximal rate $\rho_5$.
\end{remark}

% \begin{remark}[Connection to Theorem~\ref{thm:maximal-growth}]
\begin{remark}[Connection to Theorem ?????]
The quantity \(\rho_5\) defined above represents the maximal geometric growth factor
per symbol for any admissible word under the horizontal-sum dynamics.  Taking the fifth
root over one full $5$-cycle yields the maximal per-symbol growth, which directly controls 
% \(\limsup_{i\to\infty} \widetilde A(i)^{1/i}\) in Theorem~\ref{thm:maximal-growth}. 
\(\limsup_{i\to\infty} \widetilde A(i)^{1/i}\) in Theorem ????. 
Because \(\rho_5 < 3\), the growth of the number of convergent residue classes is
asymptotically strictly slower than the corresponding powers of two, guaranteeing
that the density of integers with large stopping times vanishes.
\end{remark}

\begin{remark}[Derivation and Attribution of Local Growth Constants]
The horizontal-sum framework used to analyze $\widetilde A(i)$ is inspired by
the Pascal-type summation approach introduced by Winkler~\cite{Winkler2017},
which provided the original insight that admissible Collatz division words can
be organized into structured combinatorial growth processes.  

The present work adapts this perspective to the specific setting of
\OEIS{A022921} and \OEIS{A186009}, with one key modification: unlike
Winkler's vertical-sum formulation, the horizontal orientation used here is
essential for preserving row-wise combinatorial structure, enabling precise
extremal growth calculations. This extension allows a complete analysis of
local growth configurations and the derivation of the uniform 5-cycle bound.

The constants shown in Lemma~\ref{lem:local-growth-bounds} are obtained by
examining the horizontal-sum dynamics governing $\widetilde A(i)$.
Within each local configuration, elements arising from single steps evolve
according to standard Pascal-type summation, yielding maximal growth
proportional to powers of $2$. Additional entries introduced by doubling steps
are bounded by summing on top of this baseline growth.

For each admissible local pattern occurring within \OEIS{A022921}, all possible
placements of delayed doubling contributions were examined, and the largest
achievable ratio of post-pattern to pre-pattern row sums was recorded. No
assumptions beyond the combinatorial structure of the horizontal-sum
construction are used.

A complete derivation of the horizontal-sum mechanism, together with detailed
computations of all local growth constants and their admissible concatenations,
is given in Appendix~\ref{app:local-growth-details}. Consequently, while the
conceptual framework traces back to Winkler’s original insight, the bounds
employed here are established entirely within the present paper.  Note that no
analytic or probabilistic assumptions are made; all bounds arise from finite
combinatorial enumeration.

It follows that any admissible sequence under \OEIS{A022921}, including those
relevant for \OEIS{A186009}, is strictly dominated in per-symbol growth by a
concatenation of $5$-cycles, yielding the uniform bound of
Proposition~\ref{prop:extremal-5-cycle}.
\end{remark}

% \subsection{Cycle Structure and Net Multiplicative Effect}

% A frequently occurring block is the $7$--cycle
% \[
% \{ s, d, s, d, s, d, d \},
% \]
% which contains $3$ singles and $4$ doubles.  The corresponding number of
% factors of $2$ is
% \[
% 3 + 2(4) = 11,
% \]
% so the net multiplicative contribution of this cycle satisfies
% \[
% \frac{2^{11}}{3^7} < 1.
% \]

% Thus the $7$--cycle is strictly contractive. Every $7$--cycle is followed by a $5$--cycle
% \[
% \{ s, d, s, d, d \},
% \]
% which contains $2$ singles and $3$ doubles, contributing
% \[
% 2 + 2(3) = 8
% \]
% factors of $2$.  Since
% \[
% \frac{2^8}{3^5} > 1,
% \]
% the $5$--cycle is expansive.

% Concatenating the $7$--cycle and $5$--cycle produces a $12$--cycle
% \[
% \{ s, d, s, d, s, d, d, s, d, s, d, d \},
% \]
% containing $5$ singles and $7$ doubles.  The total number of factors of $2$ is
% \[
% 5 + 2(7) = 19,
% \]
% and hence
% \[
% \frac{2^{19}}{3^{12}} < 1.
% \]
% Therefore, the $12$--cycle is again contractive.

% Repeated 12-cycles remain contractive, suppressing growth until enough
% slack accumulates to allow a 5-cycle insertion. The 5-cycle represents the
% unique locally expansive pattern, though its repetition is constrained by
% the grammar.

\subsection{Horizontal-Sum Dynamics and Pascal-Type Growth}

The generating rows evolve by a horizontal-sum rule:
\begin{equation}
a_i(n) = a_i(n-1) + a_{i-1}(n), \quad i,n \in \mathbb{N},
\label{eq:horizontal-sum-recurrence}
\end{equation}
so that sequences of single steps exhibit ordinary Pascal--type growth.
Double steps differ in that the terminal entry is duplicated, causing the
final term to appear twice while preserving monotonicity.

Because the grammar forces expansive configurations to appear only within this minimal 5-block,
extremal asymptotic growth is realized by repeated alignment of this block, making it sufficient
to analyze this configuration.  Thus for extremal 5-cycle analysis, the local maximal sequence
(dominating per-pattern growth) is
\[
\{ s, d, s, d, d \}, \{ s, d, s, d, d \},
\]
enumerating occurrences of each admissible local pattern within a 5-cycle.

\subsection*{Admissible Local Patterns}
\label{sec:admissible-patterns}

Within \OEIS{A022921}, only four local configurations occur:
\[
(s, d), \quad (d, d), \quad (s, d, s), \quad (d, d, s).
\]

The associated coefficient counts within a 5-cycle are:
\[
\begin{array}{c|c}
\text{Pattern} & \text{Associated coefficient count} \\
\hline
s,d & 2 \\
d,d & 1 \\
s,d,s & 1 \\
d,d,s & 1
\end{array}
\]

% ============================================================
\section{Beatty Structure and Tower Decomposition of the Denominator}
\label{sec:beatty-towers}
% ============================================================

The admissible grammar is controlled by the threshold sequence
\[
A020914(n)=\lfloor n\log_2 3\rfloor+1,
\]
whose increment structure is Sturmian.

Consequently, large-scale behavior of admissible division words is governed
by Diophantine approximations to \(\log_2 3\). The convergents and
semiconvergents determine distinguished interval lengths at which dyadic
and ternary growth become nearly synchronized.

We formalize the macroscopic structure governing the denominator growth
through the Beatty sequence associated with $\log_2 3$.

% ------------------------------------------------------------

\begin{definition}[Dyadic Threshold Sequence]
Let
\[
\alpha := \log_2 3,
\qquad
d(n) := A020914(n) = \lfloor n\alpha \rfloor + 1.
\]
\end{definition}

The increment word
\[
\Delta(n) := d(n+1)-d(n)
\]
takes values in $\{1,2\}$ and defines a Sturmian (Beatty) sequence.

% ------------------------------------------------------------

\begin{definition}[Tower Lengths]
A sequence of integers $\{q_k\}$ is called a tower sequence if
\[
q_k \alpha = p_k + \varepsilon_k,
\]
where $p_k \in \mathbb{Z}$ and $|\varepsilon_k|$ is minimal among all
denominators up to $q_k$.

We refer to $q_k$ as a \emph{tower length} and $p_k$ as its associated
dyadic exponent.
\end{definition}

Empirically, the tower sequence begins as
\[
5,\,7,\,12,\,41,\,53,\,306,\,359,\,665,\,15601,\,16266,\,31867,\,\dots
\]

% ------------------------------------------------------------

\begin{lemma}[Diminishing Approximation Error]
\label{lem:tower-error}
Let $\varepsilon_k := q_k \alpha - p_k$.
Then
\[
|\varepsilon_k| \to 0 \quad \text{as } k \to \infty.
\]
\end{lemma}

\begin{proof}
This follows from the theory of continued fractions and best rational
approximations of irrational numbers.
\end{proof}

% ------------------------------------------------------------

\begin{lemma}[Alternating Approximation Sign]
\label{lem:alternating}
The errors $\varepsilon_k$ alternate in sign.  That is,
tower approximations alternate between underestimates and overestimates
of $\alpha$.
\end{lemma}

\begin{proof}
This is a standard property of convergents of continued fractions.
\end{proof}

% ------------------------------------------------------------

\begin{definition}[Expansive and Contractive Towers]
A tower length $q_k$ is called:
\begin{itemize}
\item \emph{expansive} if $\varepsilon_k < 0$ (i.e.\ $q_k\alpha < p_k$),
\item \emph{contractive} if $\varepsilon_k > 0$.
\end{itemize}
\end{definition}

Thus expansive towers slightly undercount dyadic growth, while
contractive towers slightly overcount it.

% ------------------------------------------------------------

\begin{lemma}[Tower Partition Structure]
\label{lem:tower-partition}
Tower lengths occur in finite clusters of the form
\[
\{5,7,12\},\quad
\{41,53\},\quad
\{306,359,665\},\quad
\{15601,16266,31867\},\ \dots
\]
in which:
\begin{itemize}
\item the smallest element is expansive,
\item all remaining elements are contractive.
\end{itemize}
\end{lemma}

\begin{proof}[Heuristic justification]
These clusters arise from convergents and semiconvergents of $\alpha$.
The smallest element in each cluster lies below $\alpha$ (negative error),
while subsequent refinements overshoot (positive error) before the next
cycle begins.
\end{proof}

% ------------------------------------------------------------

\begin{lemma}[Exact Jump Count on Towers]
\label{lem:exact-jumps}
Let $q_k$ be a tower length.  Then the increment word over the interval
$[n,n+q_k)$ contains an exact number of double steps (jumps), independent
of $n$.
\end{lemma}

\begin{proof}[Sketch]
Since
\[
d(n+q_k)-d(n) \in \{p_k, p_k+1\},
\]
and $\varepsilon_k$ is small, the total number of increments equal to $2$
is fixed across the block.  The Sturmian structure ensures that all
length-$q_k$ factors have identical Parikh counts.
\end{proof}

% ------------------------------------------------------------

\begin{lemma}[Terminal Double--Double Structure]
\label{lem:dd-terminal}
Every tower block ends in a double--double pattern $(2,2)$.
\end{lemma}

\begin{proof}[Heuristic]
The atomic building blocks of the increment word are the $5$-cycle and
$7$-cycle, both of which terminate in $(2,2)$.  Since all tower lengths
are concatenations of these atomic cycles, the terminal structure is
preserved.
\end{proof}

% ------------------------------------------------------------

\begin{theorem}[Macroscopic Denominator Law]
\label{thm:denominator-law}
Let $q_k$ be a tower sequence.  Then over any interval of length $q_k$:
\begin{enumerate}
\item The total dyadic exponent satisfies
\[
d(n+q_k)-d(n) = p_k + O(1),
\]
with $|O(1)| \le 1$.

\item The jump count is exact and independent of $n$.

\item The multiplicative imbalance satisfies
\[
\frac{3^{q_k}}{2^{p_k}} = 2^{\varepsilon_k},
\]
where $\varepsilon_k \to 0$.

\item Expansive towers ($\varepsilon_k<0$) create slack, while
contractive towers ($\varepsilon_k>0$) consume slack.

\item The sequence of towers alternates between expansion and contraction,
with diminishing magnitude.
\end{enumerate}
\end{theorem}

% ------------------------------------------------------------

\begin{remark}[Denominator Rigidity]
The denominator is completely governed by the Beatty structure of
$\log_2 3$.  All macroscopic irregularities reduce to a controlled,
alternating sequence of small multiplicative errors across tower scales.
\end{remark}


% ============================================================
\section{Canonical Minimal Envelope and Regenerative Growth}
\label{sec:minimal-envelope}
% ============================================================

Let
\[
a(n)=A186009(n),
\]
and let
\[
d(n)=A020914(n)=\lfloor n\log_2 3\rfloor+1.
\]

Define the increment sequence
\[
\delta(n):=d(n+1)-d(n)\in\{1,2\}.
\]

A value $\delta(n)=2$ induces a terminal duplication event in the
admissible tuple grammar.

The central structural observation is that the admissible grammar is
monotone with respect to the placement of duplication events:
earlier duplication permanently increases future tuple mass.

Consequently, the canonical Beatty phase beginning at intercept $0$
produces the minimal admissible growth envelope.

% ------------------------------------------------------------

\subsection{Jump Schedules and Tuple Evolution}

% ------------------------------------------------------------

\begin{definition}[Jump Schedule]
A jump schedule is a sequence
\[
J=(\delta_0,\delta_1,\delta_2,\dots),
\qquad
\delta_i\in\{1,2\},
\]
describing the placement of terminal duplication events in the
tuple recursion.

The canonical schedule is the Beatty/Sturmian increment sequence
\[
\delta_i=d(i+1)-d(i),
\]
associated to
\[
d(i)=\lfloor i\log_2 3\rfloor+1.
\]
\end{definition}

% ------------------------------------------------------------

\begin{definition}[Generating Tuple]
Let
\[
T_n=(t_1,t_2,\dots,t_r)
\]
denote the admissible generating tuple at level $n$.

The tuple evolves by:
\begin{enumerate}
\item cumulative Pascal-type addition,
\item terminal duplication whenever $\delta(n)=2$.
\end{enumerate}

Define the total tuple mass by
\[
|T_n|:=\sum_i t_i.
\]
\end{definition}

% ------------------------------------------------------------

\subsection{Monotonicity of the Grammar}

% ------------------------------------------------------------

\begin{lemma}[Positivity Preservation]
\label{lem:positivity}
All tuple entries remain positive under admissible evolution.

Consequently, any increase introduced at one stage propagates
monotonically to all future stages.
\end{lemma}

\begin{proof}
The admissible recursion consists entirely of cumulative additions
of positive quantities.

Terminal duplication copies an existing positive terminal entry,
and no subtraction occurs anywhere in the grammar evolution.

Hence all tuple entries remain positive, and any increase introduced
at one stage propagates to all future descendants.
\end{proof}

% ------------------------------------------------------------

\begin{lemma}[Earlier Duplication Dominance]
\label{lem:earlier-duplication}
Let
\[
J_1,\qquad J_2
\]
be two admissible jump schedules that agree through some index $m$.

Suppose $J_1$ performs a terminal duplication strictly earlier than
$J_2$.

Then the corresponding tuple masses satisfy
\[
|T_n^{(1)}|
\ge
|T_n^{(2)}|
\]
for all sufficiently large $n$, with strict inequality eventually
persistent.
\end{lemma}

\begin{proof}[Heuristic mechanism]
Earlier duplication inserts additional positive mass into the tuple
recursion sooner.

By Lemma~\ref{lem:positivity}, the admissible evolution is monotone:
all future cumulative sums inherit any earlier increase.

Thus earlier duplication permanently amplifies future tuple growth,
yielding eventual domination.
\end{proof}

% ------------------------------------------------------------

\subsection{Canonical Minimal Envelope}

% ------------------------------------------------------------

\begin{theorem}[Canonical Minimal Envelope]
\label{thm:minimal-envelope}
Among all admissible phase continuations of the Sturmian jump grammar,
the canonical intercept-$0$ Beatty phase produces minimal tuple growth.

Equivalently, for every admissible phase shift $\phi$,
\[
|T_n^{(0)}|
\le
|T_n^{(\phi)}|
\]
for all sufficiently large $n$.
\end{theorem}

\begin{proof}[Structural mechanism]
The canonical Beatty phase delays terminal duplication events
as long as admissibly possible.

By Lemma~\ref{lem:earlier-duplication}, any admissible phase shift
that introduces duplication earlier produces permanently larger
future tuple mass.

Hence the canonical phase yields the minimal admissible growth
envelope.
\end{proof}

% ------------------------------------------------------------

\begin{remark}[Minimal Regenerative Orbit]
Exact self-replication of admissible tuples occurs only after
renormalization together with resetting the Beatty phase to the
canonical intercept-$0$ schedule.

However, Theorem~\ref{thm:minimal-envelope} shows that this reset
produces the smallest admissible continuation.

Thus exact regeneration is unnecessary:
the canonical reset already provides a universal lower envelope for
future admissible growth.
\end{remark}

% ============================================================
\section{Restricted Duplication Grammar and Spectral Separation}
\label{sec:restricted-grammar}
% ============================================================

The admissible tuple grammar defining \OEIS{A186009} possesses a natural
restricted subsystem obtained by forbidding consecutive terminal
duplication events.

This restricted grammar generates the sequence \OEIS{A047749}, which
therefore serves as a canonical lower-growth sector of the full
admissible grammar.

% ------------------------------------------------------------

\subsection{Restricted Duplication Grammar}

Recall that the admissible grammar evolves through:
\begin{enumerate}
\item cumulative Pascal-type propagation,
\item terminal duplication whenever the jump increment satisfies
\[
\delta(n)=2.
\]
\end{enumerate}

In the full grammar, consecutive duplication events are permitted.

The restricted grammar forbids consecutive duplication, so every
duplication event must be followed by at least one non-duplication step.

Equivalently, if
\[
D
\]
denotes a duplication step and
\[
S
\]
a non-duplication step, then the restricted grammar forbids the local
pattern
\[
DD.
\]

% ------------------------------------------------------------

\begin{definition}[Restricted Duplication Sequence]
Let
\[
a_0(n):=A047749(n).
\]

Then \(a_0(n)\) is generated by the admissible tuple recursion under
the constraint that consecutive terminal duplications are forbidden.
\end{definition}

% ------------------------------------------------------------

\subsection{Initial Coincidence and Spectral Divergence}

The restricted and unrestricted grammars initially coincide because the
canonical Beatty increment sequence contains no consecutive duplication
events at small scales.

Consequently,
\[
A047749(n)=A186009(n+1)
\]
through the first coincidence range.

However, once the unrestricted grammar encounters the first admissible
double--double event, the two systems diverge permanently.

% ------------------------------------------------------------

\begin{lemma}[Permanent Amplification by Consecutive Duplication]
\label{lem:permanent-amplification}
Suppose two admissible tuple evolutions agree through some level \(m\),
and suppose one evolution performs a consecutive terminal duplication
event at level \(m+1\) while the other does not.

Then the unrestricted evolution eventually dominates strictly:
\[
|T_n^{(\mathrm{full})}|
>
|T_n^{(\mathrm{restricted})}|
\]
for all sufficiently large \(n\).
\end{lemma}

\begin{proof}[Structural mechanism]
Terminal duplication inserts additional positive mass into the tuple.

By positivity preservation of the admissible recursion, all future
Pascal-type propagations inherit this excess mass monotonically.

Thus the effect of an earlier duplication event is permanent and
strictly amplifying.
\end{proof}

\begin{corollary}[Spectral Domination]
If the exponential growth rates exist, then
\[
\lambda_0 \le \lambda.
\]

Moreover, computational evidence and the persistent amplification
mechanism strongly suggest that the inequality is strict:
\[
\lambda_0 < \lambda.
\]
\end{corollary}

% ------------------------------------------------------------
\subsection{Interpretation}

The sequence \OEIS{A047749} therefore represents the minimal
non-consecutive amplification sector of the admissible grammar, while
\OEIS{A186009} represents the full unrestricted positive grammar.

The spectral gap between the two systems is generated entirely by the
presence of consecutive terminal duplication events.

This identifies consecutive duplication as the fundamental local
mechanism responsible for accelerated asymptotic growth of admissible
division-word configurations.

% ------------------------------------------------------------
\section{Grammar-Induced Deficit Compensation}
\label{sec:grammar-compensation}

We formalize a structural mechanism in which jumps in the dyadic
threshold sequence create temporary mass deficits, while the same jumps
simultaneously enlarge the admissible grammar of minimal descent words.
The resulting delayed compensation drives cumulative density toward $1$.

% ------------------------------------------------------------

\begin{definition}[Dyadic Threshold Sequence]
Let
\[
d(n)=A020914(n).
\]
Equivalently,
\[
d(n)=\lfloor n\log_2 3 \rfloor + 1.
\]
This is the minimal dyadic exponent required to exceed the ternary scale
$3^n$.
\end{definition}

% ------------------------------------------------------------

\begin{definition}[Jump Index]
An index $n\ge0$ is called a \emph{jump index} if
\[
d(n+1)-d(n)=2.
\]
Otherwise,
\[
d(n+1)-d(n)=1.
\]
\end{definition}

% ------------------------------------------------------------

\begin{lemma}[Binary Step Law]
\label{lem:binary-step}
For every $n\ge0$,
\[
d(n+1)-d(n)\in\{1,2\}.
\]
Hence the sequence $d(n)$ evolves only by unit steps or jump steps.
\end{lemma}

\begin{proof}
Using
\[
d(n)=\lfloor n\log_2 3\rfloor+1,
\]
we obtain
\[
d(n+1)-d(n)
=
\lfloor (n+1)\log_2 3\rfloor-\lfloor n\log_2 3\rfloor.
\]
Since
\[
1<\log_2 3<2,
\]
the floor can increase only by $1$ or $2$.
\end{proof}

% ------------------------------------------------------------

\begin{lemma}[Single Dyadic Deficit Law]
\label{lem:deficit-law}
If $n$ is a jump index, then exactly one intermediate dyadic level is
skipped relative to unit growth.

Equivalently, the denominator factor changes from the expected scale
$2^{d(n)+1}$ to the actual scale $2^{d(n)+2}$, reducing the associated
mass contribution by a factor of $1/2$.
\end{lemma}

\begin{proof}
At a jump index,
\[
d(n+1)=d(n)+2.
\]
Relative to the unit-step scale $d(n)+1$, one additional factor of $2$
appears in the denominator, so the contribution is halved.
\end{proof}

% ------------------------------------------------------------

\begin{definition}[Minimal Descent Mass]
Let
\[
N(n)=\frac{A186009(n+1)}{2^{d(n)}}.
\]
This is the total density contribution of new minimal descent words at
level $n$.
\end{definition}

% ------------------------------------------------------------

\begin{lemma}[Jump Enlarges the Admissible Grammar]
\label{lem:jump-grammar}
At every jump index, the admissible prefix-sum region defining minimal
descent words strictly enlarges at subsequent levels.

Consequently, new minimal descent words must occur after each jump.
\end{lemma}

\begin{proof}[Proof Sketch]
The admissible words are determined by prefix constraints of the form
\[
H(i)<d(i+1), \qquad H(m-1)=d(m).
\]
When $d(n)$ jumps, the terminal admissible sum increases while prior
prefix constraints remain unchanged. This creates new feasible lattice
points in the constrained composition region, yielding new admissible
words.
\end{proof}

% ------------------------------------------------------------

\begin{lemma}[Lagged Numerator Response]
\label{lem:lagged-response}
There exists an integer $L\ge1$ such that every jump at index $n$
produces at least one new minimal descent word among the levels
\[
n+1,n+2,\dots,n+L.
\]
\end{lemma}

\begin{proof}[Heuristic Basis]
Computational enumeration shows that new admissible words appear after
every jump, with uniformly bounded delay. The bound $L$ is empirically
small.
\end{proof}

% ------------------------------------------------------------

\begin{lemma}[Finite Compensation Packet]
\label{lem:packet}
There exist constants $L\ge1$ and $c>0$ such that every jump index $n$
satisfies
\[
\sum_{i=1}^{L} N(n+i)\ge c\,2^{-d(n)}.
\]
That is, each jump-induced unit deficit is compensated by a bounded
future packet of new descent mass.
\end{lemma}

% ------------------------------------------------------------

\begin{definition}[Residual Gap]
Let
\[
C(m)=\sum_{k=0}^{m}N(k),
\qquad
R(m)=1-C(m).
\]
\end{definition}

% ------------------------------------------------------------

\begin{theorem}[Compensation Implies Gap Contraction]
\label{thm:gap-contraction}
Suppose the finite compensation packet property holds uniformly.
Then there exist constants $M\ge1$ and $\epsilon>0$ such that
\[
R(m+M)\le (1-\epsilon)R(m)
\]
for all sufficiently large $m$.
\end{theorem}

\begin{proof}[Proof Sketch]
Each jump creates at most one unit deficit by
Lemma~\ref{lem:deficit-law}. By Lemma~\ref{lem:packet}, each such deficit
is repaired within bounded delay by a positive packet of future mass.
Uniform boundedness implies a fixed positive fraction of the remaining
gap is removed over each window of length $M$.
\end{proof}

% ------------------------------------------------------------

\begin{corollary}[Density Exhaustion]
\label{cor:density-exhaustion}
Under the hypotheses of Theorem~\ref{thm:gap-contraction},
\[
\lim_{m\to\infty} C(m)=1.
\]
\end{corollary}

\begin{proof}
Iterating the contraction inequality gives
\[
R(m+kM)\le (1-\epsilon)^kR(m)\to0.
\]
Hence $C(m)\to1$.
\end{proof}

% ------------------------------------------------------------

\begin{remark}[Critical Reduction]
The remaining task is purely combinatorial: prove a uniform bound on the
delay and size of compensation packets generated by jump indices in the
admissible division-word grammar.
\end{remark}


\appendix

% ============================================================
\newpage
\section{Canonical Determination of the Constant \(c(\omega)\)}
\label{app:canonical-c}
% ============================================================

Each division word \(\omega\) induces an affine map \(F_\omega(x)\) of length $m$
(see Section~\ref{sec:division-counts})
\[
F_\omega(x) = \frac{3^m x + c(\omega)}{2^{D(\omega)}}, \qquad D(\omega) = d_0 + \dots + d_{m-1},
\]
where \(c(\omega)\) is the affine constant canonically determined by the
division word \(\omega\).  The associated congruence constraint is naturally
considered modulo \(2^{D(\omega)}\).

The following procedure reconstructs the same affine constant \(c(\omega)\)
from any representative of the associated residue class.

\subsection{Canonical Choice}

For a fixed \(\omega\), any integer \(x\) in the corresponding residue class satisfies
\[
3^m x + c \equiv 0 \pmod{2^{D(\omega)}}.
\]
We remove ambiguity by selecting the unique representative
\[
0 \le c(\omega) < 2^{D(\omega)}.
\]
This choice depends only on \(\omega\) and is independent of the representative \(x\).

\subsection{Step-by-Step Computation}

Given \(\omega\) of length \(m\) and total division count \(D(\omega)\):

\begin{enumerate}
\item Choose any integer \(x\) in the residue class corresponding to \(\omega\).  
\item Compute \(y = 3^m x\).  
\item Let \(r = \lceil y / 2^{D(\omega)} \rceil\) be the smallest integer such that \(r \cdot 2^{D(\omega)} \ge y\).  
\item Define \(c(\omega) = r \cdot 2^{D(\omega)} - y\).  
\end{enumerate}

By this construction, \(0 \le c(\omega) < 2^{D(\omega)}\) ensuring that
\[
F_\omega(x) = \frac{3^m x + c(\omega)}{2^{D(\omega)}} \in \mathbb{Z} \quad \text{for all } x \text{ in the class.}
\]

\newpage
\subsection{Example}

Let \(\omega = (1,1,1,4)\) with \(m=4\) and \(D(\omega) = \sum_{i=0}^{3} d_i = 1 + 1 + 1 + 4 = 7\):

\begin{enumerate}
\item Choose \(x = 15\). Compute \(y = 3^4 \cdot 15 = 1215\).  
\item \(2^7 = 128\), so \(r = \lceil 1215/128 \rceil = 10\).  
\item \(c(\omega) = 10 \cdot 128 - 1215 = 65\).  
\end{enumerate}

Hence the canonical map is
\[
F_\omega(x) = \frac{81x + 65}{128}.
\]

Checking another representative \(x = 143\):
\[
F_\omega(143) = \frac{81 \cdot 143 + 65}{128} = 91,
\]
confirming that the same \(c(\omega)\) works for the entire residue class.

\subsection*{Remark}

This construction is purely illustrative and is not required for the density
or growth arguments in the main text. It demonstrates concretely how the affine
map is fully determined by the division word \(\omega\).

Since any two integers realizing the same division word differ by a multiple of $2^{D(\omega)}$,
the value of $c(\omega)$ computed by this procedure is independent of the chosen
representative.

\begin{remark}[Cylinder Rigidity]
The affine constant \(c(\omega)\) is invariant across the entire residue
class realizing \(\omega\).  Consequently, the symbolic data encoded by
the division word uniquely determines both the affine map and the associated
cylinder modulo \(2^{D(\omega)}\).
\end{remark}


% ============================================================
\newpage
\section{Cylinder Reconstruction and Affine Structure}
\label{app:cylinder-reconstruction}
% ============================================================

This appendix illustrates how admissible division words determine
explicit congruence cylinders and corresponding affine arithmetic
families.

% ------------------------------------------------------------

\subsection{Cylinder Reconstruction}

Let
\[
\omega=(d_0,\dots,d_{m-1})
\]
be an admissible division word with total division count
\[
D(\omega)=\sum_{i=0}^{m-1} d_i.
\]

From the affine representation,
\[
F_\omega(x)
=
\frac{3^m x+c(\omega)}{2^{D(\omega)}}.
\]

Admissibility is equivalent to the congruence condition
\[
3^m x+c(\omega)\equiv0\pmod{2^{D(\omega)}}.
\]

Since \(3^m\) is invertible modulo \(2^{D(\omega)}\), there exists a
unique residue class
\[
x\equiv r_\omega \pmod{2^{D(\omega)}}.
\]

Thus all integers realizing \(\omega\) are parametrized by
\[
x(k)=r_\omega+k\,2^{D(\omega)},
\qquad
k\in\mathbb N_0.
\]

The corresponding congruence class is called the
\emph{cylinder} associated with \(\omega\).

% ------------------------------------------------------------

\subsection{Affine Arithmetic Structure}

Applying the affine map to \(x(k)\) gives
\[
F_\omega(x(k))
=
\frac{3^m(r_\omega+k2^{D(\omega)})+c(\omega)}
     {2^{D(\omega)}}.
\]

Expanding,
\[
F_\omega(x(k))
=
k\,3^m
+
\frac{3^m r_\omega+c(\omega)}
     {2^{D(\omega)}}.
\]

Define
\[
K_\omega
:=
\frac{3^m r_\omega+c(\omega)}
     {2^{D(\omega)}}.
\]

Then
\[
F_\omega(x(k))
=
k\,3^m+K_\omega.
\]

Thus every admissible division word induces:

\begin{itemize}
\item a congruence cylinder modulo \(2^{D(\omega)}\), and
\item an affine arithmetic progression with slope \(3^m\).
\end{itemize}

Consequently, the symbolic dynamics partitions the odd integers into
affine congruence fibers determined entirely by the division grammar.

% ------------------------------------------------------------

\subsection{Example}

Consider the admissible division word
\[
\omega=(1,1,1,4),
\]
for which
\[
m=4,
\qquad
D(\omega)=7,
\]
and
\[
F_\omega(x)=\frac{81x+65}{128}.
\]

The admissibility condition becomes
\[
81x+65\equiv0\pmod{128}.
\]

Solving yields the unique residue class
\[
x\equiv15\pmod{128}.
\]

Hence the cylinder is
\[
x(k)=15+128k.
\]

Applying the affine map,
\[
F_\omega(15+128k)
=
81k+10.
\]

Thus the symbolic cylinder generated by \(\omega\) is mapped onto the
arithmetic progression
\[
10,\,91,\,172,\,253,\dots
\]
with common difference \(81=3^4\).

% ------------------------------------------------------------

\subsection{Illustrative Derivation of the Affine Constant}

For readers wishing to see the affine constant arise directly from
iterated substitution, consider a general division word
\[
\omega=(a,b,c,d)
\]
of length \(4\).

The accelerated relations are
\[
x_1=\frac{3x_0+1}{2^a},
\qquad
x_2=\frac{3x_1+1}{2^b},
\qquad
x_3=\frac{3x_2+1}{2^c},
\qquad
x_4=\frac{3x_3+1}{2^d}.
\]

Eliminating intermediate variables recursively yields
\[
2^{a+b+c+d}x_4
=
3^4x_0
+
3^3
+
3^2 2^a
+
3\,2^{a+b}
+
2^{a+b+c}.
\]

Hence
\[
F_\omega(x)
=
\frac{
3^4x
+
\left(
3^3
+
3^2 2^a
+
3\,2^{a+b}
+
2^{a+b+c}
\right)
}{
2^{a+b+c+d}
},
\]
so that
\[
c(\omega)
=
3^3
+
3^2 2^a
+
3\,2^{a+b}
+
2^{a+b+c}.
\]

This illustrates concretely how the affine constant accumulates the
propagated contributions of successive ``+1'' terms throughout the
accelerated iteration.

% ============================================================
\newpage
\section{Continued Fractions, Tower Cycles, and the Extremal Envelope}
\label{app:continued-fractions}
% ============================================================

This appendix records the Diophantine structure underlying the admissible
grammar and explains the origin of the distinguished cycle lengths
\[
1,\ 1,\ 2,\ 5,\ 12,\ 41,\ 53,\ 306,\ 665,\ 15601,\ 31867,\ 79335,\ldots
\]
that appear throughout the paper.

% \subsection{Continued-Fraction Convergents of \texorpdfstring{$\log_2 3$}{log2(3)}}

% The continued-fraction expansion of
% \[
% \alpha=\log_2 3
% \]
% begins
% \[
% \alpha=[1;1,1,2,2,3,1,5,2,23,2,2,\dots].
% \]

% The corresponding convergents are
% \[
% \frac{1}{1},\ 
% \frac{2}{1},\ 
% \frac{3}{2},\ 
% \frac{8}{5},\ 
% \frac{19}{12},\ 
% \frac{65}{41},\ 
% \frac{84}{53},\ 
% \frac{485}{306},\ 
% \frac{1054}{665},\ 
% \frac{24727}{15601},\ 
% \frac{50508}{31867},\ 
% \frac{125743}{79335},
% \ldots
% \]

% Let
% \[
% \frac{p_k}{q_k}
% \]
% denote the $k^{th}$ convergent.  The denominators $q_k$ are
% \[
% 1,\ 1,\ 2,\ 5,\ 12,\ 41,\ 53,\ 306,\ 665,\ 15601,\ 31867,\ 79335,\ldots
% \]
% are precisely the distinguished tower lengths appearing throughout the admissible grammar.

% \subsection{Recursive Generation of Tower Cycles}

% A fundamental property of continued fractions is the recurrence
% \[
% q_k=a_k q_{k-1}+q_{k-2}.
% \]
% where $a_{k}$ is the corresponding continued-fraction coefficient.

% \[
% \begin{array}{r|r|r|r}
% k & a_k & p_k & q_k \\
% \hline
%  0 &  - & 1         & 1 \\
%  1 &  1 & 2         & 1 \\
%  2 &  1 & 3         & 2 \\
%  3 &  2 & 8         & 5 \\
%  4 &  2 & 19        & 12 \\
%  5 &  3 & 65        & 41 \\
%  6 &  1 & 84        & 53 \\
%  7 &  5 & 485       & 306 \\
%  8 &  2 & 1054      & 665 \\
%  9 & 23 & 24727     & 15601 \\
% 10 &  2 & 50508     & 31867 \\
% 11 &  2 & 125743    & 79335
% \end{array}
% \]

% Consequently every tower length is assembled canonically from the two preceding tower
% lengths.  Defining atomic cycles $C_0 = 1$ as a contractive single exponent of 2
% ($\frac{2}{3}$) and $C_1 = 2$ as an expansive double exponent of 2 ($\frac{4}{3}$),
% the cycle hierarchy is generated using the $q_k=a_k q_{k-1}+q_{k-2}$ recursive assembly rule.
\subsection{Continued-Fraction Origin of the Tower Cycles}

Let

\[
\alpha=\log_2 3.
\]

Its continued-fraction expansion begins

\[
\alpha=[1;1,1,2,2,3,1,5,2,23,2,2,\dots].
\]

Writing

\[
\frac{p_k}{q_k}
\]

for the convergents of $\alpha$, the first denominator sequence is

\[
1,\ 1,\ 2,\ 5,\ 12,\ 41,\ 53,\ 306,\ 665,\ 15601,\ldots
\]

which coincides exactly with the tower lengths observed in the admissible grammar.
The convergent denominators satisfy the standard recurrence

\[
q_k=a_k q_{k-1}+q_{k-2},
\]

where $a_k$ is the corresponding continued-fraction coefficient.
Consequently every tower cycle is assembled canonically from the two
preceding tower cycles.

Table~\ref{tab:tower-cycles} records the first instances.

\begin{table}[ht]
\centering
\caption{Continued-fraction generation of the tower cycles.}
\label{tab:tower-cycles}
\begin{tabular}{r l r r l l}
\hline
$q_k$ & Cycle & $k$ & $a_k$ & Cycle Composition & Sequence \\
\hline
1     & $C_{    0}$ &  0 & -- & Single                          & $1$ \\
1     & $C_{    1}$ &  1 &  1 & Double                          & $2$ \\
2     & $C_{    2}$ &  2 &  1 & $C_{    0} + 1 \cdot C_{    1}$ & $1,2$ \\
5     & $C_{    5}$ &  3 &  2 & $2 \cdot C_{    2} + C_{    1}$ & $1,2,1,2,2$ \\
12    & $C_{   12}$ &  4 &  2 & $C_{    2} + 2 \cdot C_{    5}$ & $1,2,1,2,1,2,2,1,2,1,2,2$ \\
41    & $C_{   41}$ &  5 &  3 & $3 \cdot C_{   12} + C_{    5}$ & $1,2,1,2,1,2,2,\ldots,1,2,1,2,2$  \\
53    & $C_{   53}$ &  6 &  1 & $C_{   12} + 1 \cdot C_{   41}$ & $1,2,1,2,1,2,2,\ldots,1,2,1,2,2$  \\
306   & $C_{  306}$ &  7 &  5 & $5 \cdot C_{   53} + C_{   41}$ & $1,2,1,2,1,2,2,\ldots,1,2,1,2,2$  \\
665   & $C_{  665}$ &  8 &  2 & $C_{   53} + 2 \cdot C_{  306}$ & $1,2,1,2,1,2,2,\ldots,1,2,1,2,2$  \\
15601 & $C_{15601}$ &  9 & 23 & $23 \cdot C_{665} + C_{   306}$ & $1,2,1,2,1,2,2,\ldots,1,2,1,2,2$  \\
31867 & $C_{31867}$ & 10 &  2 & $C_{  665} + 2 \cdot C_{15601}$ & $1,2,1,2,1,2,2,\ldots,1,2,1,2,2$  \\
79335 & $C_{79335}$ & 11 &  2 & $2 \cdot C_{31867} + C_{15601}$ & $1,2,1,2,1,2,2,\ldots,1,2,1,2,2$  \\
\hline
\end{tabular}
\end{table}

% Commencing with the first concatentation of atomic sequences, namely $C_2$, each the tower
% cycle begins with a contractive sequence and terminates in an expansive sequence, where the
% internal structure is determined by the continued-fraction coefficients of $\log_2 3$.

% Thus every longer admissible cycle is obtained by concatenating previous cycles
% according to the continued-fraction coefficients of $\log_2 3$.
% The tower hierarchy is therefore not an empirical feature of the grammar but a
% direct consequence of the Diophantine approximation structure of $\log_2 3$.

% \begin{remark}[Canonical nature of the tower hierarchy]
% The cycle lengths
% \[
% 1,\ 1,\ 2,\ 5,\ 12,\ 41,\ 53,\ 306,\ 665,\ 15601,\ldots
% \]
% are not selected by observation or computation. They arise uniquely as the
% denominators of the convergents of $\log_2 3$ and therefore represent the
% canonical scales at which dyadic and ternary growth are optimally synchronized.
% The alternating expansive and contractive behavior observed in the admissible
% grammar is a direct manifestation of this Diophantine structure.
% \end{remark}

% Thus every higher tower length is canonically assembled from previous tower
% lengths. The cycle hierarchy is therefore not an empirical feature of the
% grammar; it is forced by the continued-fraction expansion of $\log_2 3$.

Thus the tower hierarchy is not an empirical feature of the grammar.
It is forced by the Diophantine approximation structure of $\log_2 3$.

\subsection{Alternating Expansion and Contraction}

The convergents alternately overestimate and underestimate $\alpha$.

Define
\[
\varepsilon_k = q_k\alpha-p_k.
\]

Then
\[
\varepsilon_k \to 0
\]
and the signs alternate.

Since
\[
\alpha=\log_2 3,
\]
it follows that
\[
\varepsilon_k = q_k\log_2 3-p_k = -\log_2\!\left(\frac{2^{p_k}}{3^{q_k}}\right).
\]

Equivalently,
\[
\frac{2^{p_k}}{3^{q_k}} = 2^{-\varepsilon_k}.
\]

% Consequently,
% \[
% \left|
% \log\!\left(\frac{2^{p_k}}{3^{q_k}}\right)
% \right|
% \to 0,
% \]
% and therefore
% \[
% \frac{2^{p_k}}{3^{q_k}}
% \to 1.
% \]
Hence

\[
\left|
\log\!\left(
\frac{2^{p_k}}{3^{q_k}}
\right)
\right|
=
(\log 2)\,|\varepsilon_k|
\to0.
\]

Therefore

\[
\frac{2^{p_k}}{3^{q_k}}
\to1.
\]

The convergents alternately overestimate and underestimate
$\log_2 3$, so the ratios alternately lie above and below unity while
approaching exact multiplicative balance.

% The ratios alternate between values greater than and less than $1$ while
% converging monotonically toward exact multiplicative balance. Examples include:

% \[
% \frac{2^8}{3^5}
% =
% \frac{256}{243}
% =
% 1.05349\ldots
% >
% 1,
% \]

% \[
% \frac{2^{19}}{3^{12}}
% =
% \frac{524288}{531441}
% =
% 0.9865\ldots
% <
% 1,
% \]

% \[
% \frac{2^{65}}{3^{41}}
% =
% 1.0115\ldots
% >
% 1,
% \]

% \[
% \frac{2^{84}}{3^{53}}
% =
% 0.9979\ldots
% <
% 1.
% \]

\begin{table}[ht]
\centering
\caption{Convergents of $\log_2 3$ and resulting
dyadic--ternary synchronization ratios.}
\label{tab:tower-ratios}
\begin{tabular}{r r c c}
\hline
$q_k$ & $p_k$ &
$\varepsilon_k=q_k\log_2(3)-p_k$ &
$\displaystyle 2^{p_k} / 3^{q_k}$ \\
\hline
 1 &  2 & $-0.41594\ldots$ & $1.33333\ldots$ \\
 2 &  3 & $+0.16993\ldots$ & $0.88889\ldots$ \\
 5 &  8 & $-0.07519\ldots$ & $1.05350\ldots$ \\
12 & 19 & $+0.01955\ldots$ & $0.98654\ldots$ \\
41 & 65 & $-0.01654\ldots$ & $1.01153\ldots$ \\
53 & 84 & $+0.00301\ldots$ & $0.99791\ldots$ \\
\hline
\end{tabular}
\end{table}

% Thus the distinguished tower lengths are precisely the scales at which
% powers of $2$ and powers of $3$ achieve exceptionally close multiplicative
% synchronization. The convergent denominators of $\log_2 3$ therefore mark
% successively better approximations to the balance point
% \[
% 2^{p_k}\approx 3^{q_k},
% \]
% explaining the alternating expansive and contractive behavior observed
% throughout the admissible grammar.

% \subsection{Return Words and Primitive Cycles}
\subsection{Primitive Return Words and the Failure of the Double Process}

The Sturmian word generated by
\[
d(n)=\lfloor n\alpha\rfloor+1
\]
admits a return-word decomposition.  At the first nontrivial level, the primitive
return words have lengths $2$, $5$, and $12$. These correspond respectively to:

\[
C_2=(1,2).
\]
\[
C_5=(1,2,1,2,2).
\]
\[
C_{12}=(1,2,1,2,1,2,2,1,2,1,2,2).
\]

The 2-cycle is locally contractive, while the 5-cycle is locally expansive. A 12-cycle
is formed by concatenating the 2-cycle with a pair of consecutive 5-cycles, but the 
12-cycle is overall contractive.

% All larger admissible structures arise from concatenations dictated by the
% continued-fraction hierarchy above with 5-cycles local expansive components.

% \subsection{Why the 2-Cycle Is Not the Extremal Envelope}

One might ask why the extremal comparison process is based on the
5-cycle rather than the shorter expansive 1-cycle corresponding to a
double increment.

Indeed, the shortest expansive block is

\[
C_1=(2),
\]

for which

\[
\frac{2^2}{3} = \frac43 > 1.
\]

% However, repeated doubles generate a much larger growth process under the
% horizontal-sum dynamics. Starting from a single entry, repeated duplication
% of the terminal term produces the rows

% \[
% 1,
% \]
% \[
% 1,\;1,
% \]
% \[
% 1,\;2,\;2,
% \]
% \[
% 1,\;3,\;5,\;5,
% \]
% \[
% 1,\;4,\;9,\;14,\;14,
% \]
% \[
% 1,\;5,\;14,\;28,\;42,\;42,
% \]
% \[
% 1,\;6,\;20,\;48,\;90,\;132,\;132,
% \]
% whose row sums are
% \[
% 1,\;2,\;5,\;14,\;42,\;132,\;429,\ldots
% \]

% These are the Catalan numbers \OEIS{A000108}.
% Consequently the consecutive growth ratios are

% \[
% \frac21,
% \qquad
% \frac52,
% \qquad
% \frac{14}{5},
% \qquad
% \frac{42}{14}=3,
% \qquad
% \frac{132}{42},
% \qquad
% \frac{429}{132},
% \ldots
% \]

Repeated doubles generate the row sums

\[
1,2,5,14,42,132,429,\ldots
\]

which are the Catalan numbers.
Since

% \[
% C_n\sim
% \frac{4^n}{n^{3/2}\sqrt{\pi}},
% \]

% their exponential growth rate equals \(4\).

% Consequently an infinite concatenation of doubles grows faster than the
% critical value \(3\) and cannot be used as an extremal comparison
% process.

% and converge to 4. Indeed, the Catalan asymptotic formula

\[
C_n
\sim
\frac{4^n}{n^{3/2}\sqrt{\pi}}
\]

implies

\[
\lim_{n\to\infty}
\frac{C_{n+1}}{C_n}
=
4.
\]

Thus an infinite concatenation of doubles produces exponential growth
rate \(4\), which exceeds the critical value \(3\).

The pure-double process is therefore unsuitable as a comparison
envelope: it overestimates growth by too large a margin and cannot
yield the desired spectral bound.

% The next shortest expansive structure is the primitive return word
% \[
% C_5=(1,2,1,2,2).
% \]

% Unlike the double process, repeated copies of \(C_5\) admit a direct
% evaluation of the associated cocycle growth and yield a per-symbol rate
% strictly below \(3\).

% For this reason \(C_5\) is the shortest expansive primitive return word
% capable of serving as a useful extremal comparison process.

\subsection{The Extremal Envelope}

Define the multiplicative cocycle
\[
G(w)
=
\frac{2^{\#s(w)+2\#d(w)}}{3^{|w|}}.
\]

where \(\#s(w)\) and \(\#d(w)\) count the number of single and double steps
in the word \(w\), respectively.

By the return-word decomposition theorem, every admissible infinite word
can be written uniquely as a concatenation of primitive return words.

Among these primitive return words, the 5-cycle
\[
C_5=(1,2,1,2,2)=(s,d,s,d,d)
\]
has maximal normalized cocycle growth. Consequently,
for every admissible return word $R$,
\[
G(R)^{1/|R|}
\le
G(C_5)^{1/5}.
\]

Although the periodic word
\[
C_5^\infty
\]
is not itself admissible in the Sturmian language, replacing each return
word in an admissible decomposition by $C_5$ can only increase the
per-symbol growth rate. Thus $C_5^\infty$ provides a universal extremal
comparison process.

\begin{proposition}[Extremal Envelope]
Let $\lambda$ denote the exponential growth rate of admissible
configurations. Then
\[
\lambda
\le
\lambda_{\mathrm{env}},
\]
where
\[
\lambda_{\mathrm{env}}
=
G(C_5)^{1/5}.
\]
\end{proposition}

\begin{proof}
Every admissible word admits a decomposition into primitive return words.
By cocycle multiplicativity, the logarithmic growth rate is a weighted
average of the normalized cocycle contributions of these return words.

Since $C_5$ has maximal normalized cocycle growth, no admissible
decomposition can exceed the per-symbol growth rate obtained by replacing
every return word with $C_5$. Therefore
\[
\lambda
\le
G(C_5)^{1/5}.
\]
\end{proof}

The extremal envelope is deliberately larger than the admissible growth
process. Indeed, the comparison word $C_5^\infty$ is forbidden by the
grammar. Consequently, if even this exaggerated comparison process
satisfies
\[
G(C_5)^{1/5}<3,
\]
then every admissible growth process must also satisfy the same bound.

\subsection{Interpretation}

The tower hierarchy, the alternating expansive and contractive blocks,
and the extremal-envelope construction all originate from the same
Diophantine source: the continued-fraction approximation of
\(\log_2 3\).

The admissible grammar therefore does not generate the distinguished
cycle lengths independently. Rather, it realizes combinatorially the
canonical convergent structure of \(\log_2 3\).

From this perspective the 5-cycle appears naturally as the first
nontrivial expansive return word, while the larger tower cycles
\(12,41,53,306,\ldots\) arise recursively through the convergent
recurrence. The extremal-envelope bound is therefore ultimately a
consequence of the Diophantine structure of \(\log_2 3\).

% ============================================================
\newpage
\section{Extremal Local Growth Calculations}
\label{app:local-growth-details}
% ============================================================

This appendix provides the explicit finite computations underlying the local
growth bounds used in the density argument. Recall that the sequence \OEIS{A020914}
has first differences given by \OEIS{A022921}, consisting of repeated blocks of
\emph{single} steps ($s$) and \emph{double} steps ($d$).


% \subsection{Cycle Structure and Net Multiplicative Effect}

% A frequently occurring block is the $7$--cycle
% \[
% \{ s, d, s, d, s, d, d \},
% \]
% which contains $3$ singles and $4$ doubles.  The corresponding number of
% factors of $2$ is
% \[
% 3 + 2(4) = 11,
% \]
% so the net multiplicative contribution of this cycle satisfies
% \[
% \frac{2^{11}}{3^7} < 1.
% \]

% Thus the $7$--cycle is strictly contractive. Every $7$--cycle is followed by a $5$--cycle
% \[
% \{ s, d, s, d, d \},
% \]
% which contains $2$ singles and $3$ doubles, contributing
% \[
% 2 + 2(3) = 8
% \]
% factors of $2$.  Since
% \[
% \frac{2^8}{3^5} > 1,
% \]
% the $5$--cycle is expansive.

% Concatenating the $7$--cycle and $5$--cycle produces a $12$--cycle
% \[
% \{ s, d, s, d, s, d, d, s, d, s, d, d \},
% \]
% containing $5$ singles and $7$ doubles.  The total number of factors of $2$ is
% \[
% 5 + 2(7) = 19,
% \]
% and hence
% \[
% \frac{2^{19}}{3^{12}} < 1.
% \]
% Therefore, the $12$--cycle is again contractive.

% Repeated 12-cycles remain contractive, suppressing growth until enough
% slack accumulates to allow a 5-cycle insertion. The 5-cycle represents the
% unique locally expansive pattern, though its repetition is constrained by
% the grammar.

% To visualize the interaction between contractive and expansive cycles,
% Figure~\ref{fig:cycle-highlight-lower} depicts the state transitions of
% \OEIS{A022921}. The red edges indicate the locally expansive 5-cycle, the
% blue edges the contractive 7-cycle, and purple edges highlight shared transitions.

% \begin{figure}[ht]
%     \centering
%     \begin{tikzpicture}
%         % Bounding box
%         \useasboundingbox (0,0) -- (10.5,6.5);

%         % Nodes
%         \node[draw=black, line width=4, circle, minimum size=2cm] (ss) at (1,5) {\Large \textbf{s, s}};
%         \node[draw=black, line width=2, circle, minimum size=2cm] (sd1) at (5,5) {\Large \textbf{s, d}};
%         \node[draw=black, line width=2, circle, minimum size=2cm] (sd2) at (9,5) {\Large \textbf{s, d}};
%         \node[draw=black, line width=2, circle, minimum size=2cm] (sd3) at (9,1) {\Large \textbf{s, d}};
%         \node[draw=black, line width=2, circle, minimum size=2cm] (d)   at (5,1) {\Large \textbf{d}};

%         % Edges
%         % Contractive 7-cycle: blue
%         \draw[->, line width=2, blue] (9,4) -- (9,2);        % sd2 -> sd3
%         \draw[->, line width=2, blue] (8,1) -- (6,1);        % sd3 -> d

%         % Expansive 5-cycle: red
%         \draw[->, line width=2, red] (8.3,4.3) -- (5.7,1.7); % sd3 -> ss

%         % Shared edge (between 5 and 7): purple
%         \draw[->, line width=2, purple] (2,5) -- (4,5);       % ss -> sd1
%         \draw[->, line width=2, purple] (5,2) -- (5,4);       % d -> sd1
%         \draw[->, line width=2, purple] (6,5) -- (8,5);       % sd1 -> sd2

%         % Legend
%         \node at (11,6)   {\color{purple} shared edge};
%         \node at (11,5.5) {\color{red}    5-cycle};
%         \node at (11,5)   {\color{blue}   7-cycle};

%     \end{tikzpicture}
%     \caption{State Transition Diagram for A022921 highlighting cycles: contractive 7-cycle (blue),
%     locally expansive 5-cycle (red), shared edges (purple). Lowercase letters reflect single ($s$)
%     and double ($d$) division steps consistent with the main text notation.}
%     \label{fig:cycle-highlight-lower}
% \end{figure}

% The figure highlights the locally expansive 5-cycle (red), the contractive 7-cycle (blue), 
% and the shared edges (purple), visually illustrating the grammar constraints of the 
% sequence \OEIS{A022921}.

% \newpage
% \subsection*{Horizontal-Sum Dynamics and Pascal-Type Growth}

% The generating rows evolve by a horizontal-sum rule:
% \begin{equation}
% a_i(n) = a_i(n-1) + a_{i-1}(n), \quad i,n \in \mathbb{N},
% % \label{eq:horizontal-sum-recurrence}
% \end{equation}
% so that sequences of single steps exhibit ordinary Pascal--type growth.
% Double steps differ in that the terminal entry is duplicated, causing the
% final term to appear twice while preserving monotonicity.

% Because the grammar forces expansive configurations to appear only within this minimal 5-block,
% extremal asymptotic growth is realized by repeated alignment of this block, making it sufficient
% to analyze this configuration.  Thus for extremal 5-cycle analysis, the local maximal sequence
% (dominating per-pattern growth) is
% \[
% \{ s, d, s, d, d \}, \{ s, d, s, d, d \},
% \]
% enumerating occurrences of each admissible local pattern within a 5-cycle.

% \subsection*{Admissible Local Patterns}
% % \label{sec:admissible-patterns}

% Within \OEIS{A022921}, only four local configurations occur:
% \[
% (s, d), \quad (d, d), \quad (s, d, s), \quad (d, d, s).
% \]

% The associated coefficient counts within a 5-cycle are:
% \[
% \begin{array}{c|c}
% \text{Pattern} & \text{Associated coefficient count} \\
% \hline
% s,d & 2 \\
% d,d & 1 \\
% s,d,s & 1 \\
% d,d,s & 1
% \end{array}
% \]

\subsection{Illustrative Extremal Computations}

The generating rows evolve by a horizontal--sum rule (see Equation \ref{eq:horizontal-sum-recurrence})
so that sequences of single steps exhibit ordinary Pascal--type growth.
Double steps differ only in that the terminal entry of the row is duplicated,
causing the final term to appear twice in the generating tuple while preserving
monotonicity of the preceding entries.


\begin{lemma}[Pascal-Type Horizontal Growth]
\label{lem:pascal-growth}

% Let a row generated by the horizontal-sum construction have total sum $S$.
Let a row generated by the horizontal-sum construction have strictly positive entries
with total sum $S$.

After one horizontal-sum step without duplication, the resulting row has total

\[
S_{\mathrm{new}}
=
S+S_{\mathrm{int}},
\]

where $S_{\mathrm{int}}$ denotes the sum of the nonterminal entries of the
previous row.

Consequently

\[
S_{\mathrm{new}} < 2S.
\]

\end{lemma}

\begin{proof}

The terminal entry of the new row equals the total sum $S$ of the previous row.
Every nonterminal entry contributes exactly once to the new row through the
horizontal-sum recurrence. Hence

\[
S_{\mathrm{new}}
=
S+S_{\mathrm{int}}.
\]

Since all entries are strictly positive,

\[
S_{\mathrm{int}}<S,
\]

and therefore

\[
S_{\mathrm{new}}<2S.
\]

\end{proof}

\begin{corollary}[Controlled Growth Under Duplication]
\label{lem:controlled-growth-under-duplication}

A subsequent duplication contributes at most one additional copy of the
terminal entry, whose value is $S$.

Hence

\[
S_{\mathrm{new}}
<
2S+S
=
3S.
\]

\end{corollary}



% \begin{lemma}[Controlled Growth of Horizontal Sums]
% \label{lem:controlled-horizontal-growth}

% Let a row have strictly positive entries with total sum $S$.

% A single increment step produces a new row whose total sum is $S$.

% If a double increment is applied immediately afterward, producing a row in which
% the terminal entry equals the sum of the previous row and all other entries remain
% strictly positive, then the total sum of the resulting row satisfies
% \[
% S_{\mathrm{new}} < 3S.
% \]
% \end{lemma}

% \begin{proof}
% The total sum of the original row is $S$.

% Under a double increment step, the resulting row consists of:
% \begin{enumerate}
%     \item (i) interior entries with total $S_{\mathrm{int}} < S$, and
%     \item (ii) two terminal contributions, each equal to $S$.
% \end{enumerate}

% Hence the total sum satisfies
% \[
% S_{\mathrm{new}} = S_{\mathrm{int}} + 2S < 3S.
% \]
% \end{proof}

These examples are chosen so that each step saturates the inequalities of
Lemma~\ref{lem:pascal-growth}, thereby realizing extremal growth. Equalities
shown below should be interpreted as extremal upper-bound realizations rather than
identities valid for all rows.

The following row-by-row computations realize extremal growth for the admissible
local patterns and together determine the sharpest uniform bounds used in
Lemma~\ref{lem:local-growth-bounds}.  Each example illustrates how a local pattern
can saturate its maximal growth factor under the horizontal-sum dynamics.  Let $S$
denote the total sum of entries in the generating row prior to a given local pattern.

\paragraph{Example 1.}
\[
\begin{array}{l}
\textit{Single } a(10) = \underbrace{173}_{\text{sum} := S}:
\underbrace{1,7,25,55,85}_{S}
\end{array}
\]

\[
\begin{array}{l}
\textit{Double } a(11) = \underbrace{476}_{\text{sum} = 3S}:
\underbrace{1,8,33,88}_{S}
\underbrace{173}_{S},
\underbrace{173}_{S}
\end{array}
\]

\[
\begin{array}{l}
\textit{Single } a(12) = \underbrace{961}_{\text{sum} = 9S}:
\underbrace{1,9,42,130}_{2S}
\underbrace{303}_{3S},
\underbrace{476}_{4S}
\end{array}
\]

\[
\begin{array}{l}
\textit{Double } a(13) = \underbrace{2652}_{\text{sum} = 33S}:
\underbrace{1,10,52,182}_{4S},
\underbrace{485}_{7S},
\underbrace{961}_{11S},
\underbrace{961}_{11S}
\end{array}
\]

\[
\begin{array}{l}
\textit{Double } a(14) = \underbrace{8045}_{\text{sum} = 123S}:
\underbrace{1,11,63,245}_{8S},
\underbrace{730}_{15S},
\underbrace{1691}_{26S},
\underbrace{2652}_{37S},
\underbrace{2652}_{37S}
\end{array}
\]

\[
\begin{array}{l}
\textit{Single } a(15) = \underbrace{17637}_{\text{sum} = 329S}:
\underbrace{1,12,75,320}_{16S},
\underbrace{1050}_{31S},
\underbrace{2741}_{57S},
\underbrace{5393}_{94S},
\underbrace{8045}_{131S}
\end{array}
\]

The resulting maximal local ratios are
\[
\begin{array}{c|c}
\text{Pattern} & \text{Max Ratio} \\
\hline
s,d & 3/1 \\
d,d & 123/33 \\
s,d,s & 9/3 \\
d,d,s & 329/123
\end{array}
\]

\paragraph{Example 2.}
\[
\begin{array}{l}
\textit{Single } a(12) = \underbrace{961}_{\text{sum} := S}:
\underbrace{1,9,42,130,303,476}_{S}
\end{array}
\]

\[
\begin{array}{l}
\textit{Double } a(13) = \underbrace{2652}_{\text{sum} = 3S}:
\underbrace{1,10,52,182,485}_{S},
\underbrace{961}_{S},
\underbrace{961}_{S}
\end{array}
\]

\[
\begin{array}{l}
\textit{Double } a(14) = \underbrace{8045}_{\text{sum} = 13S}:
\underbrace{1,11,63,245,730}_{2S},
\underbrace{1691}_{3S},
\underbrace{2652}_{4S},
\underbrace{2652}_{4S}
\end{array}
\]

\[
\begin{array}{l}
\textit{Single } a(15) = \underbrace{17637}_{\text{sum} = 37S}:
\underbrace{1,12,75,320,1050}_{4S},
\underbrace{2741}_{7S},
\underbrace{5393}_{11S},
\underbrace{8045}_{15S}
\end{array}
\]

\[
\begin{array}{l}
\textit{Double } a(16) = \underbrace{51033}_{\text{sum} = 131S}:
\underbrace{1,13,88,408,1458}_{8S},
\underbrace{4199}_{15S},
\underbrace{9592}_{26S},
\underbrace{17637}_{41S},
\underbrace{17637}_{41S}
\end{array}
\]

\[
\begin{array}{l}
\textit{Single } a(17) = \underbrace{108950}_{\text{sum} = 341S}:
\underbrace{1,14,102,510,1968}_{16S},
\underbrace{6167}_{31S},
\underbrace{15759}_{57S},
\underbrace{33396}_{98S},
\underbrace{51033}_{139S}
\end{array}
\]

The corresponding maximal local ratios are
\[
\begin{array}{c|c}
\text{Pattern} & \text{Max Ratio} \\
\hline
s,d & 3/1 \\
d,d & 13/3 \\
s,d,s & 341/131 \\
d,d,s & 37/13
\end{array}
\]

\subsection*{Improved Extremal Alignment}

Because the total multiplicative growth over a cycle is the product of the local
growth factors, the geometric mean over one full cycle is maximized by selecting
the maximal admissible factor at each local pattern occurrence. Consequently,
once a uniform upper bound has been established for each admissible local pattern,
the extremal global growth rate over a complete cycle is obtained by multiplying
these maximal local factors according to their frequencies and taking the
corresponding geometric mean.

\emph{Equivalently, because total growth over any admissible word is multiplicative in
its local factors, the asymptotic per-symbol growth rate is the geometric mean of these
factors;  hence the maximal rate over the entire language generated by the grammar is
achieved by the word that maximizes this mean subject to the grammar constraints.}

Each example yields a valid (not necessarily sharp) upper bound on the maximal growth
factor associated with each admissible local pattern.  Since these bounds arise from
the same uniform recurrence and apply globally to all admissible configurations, the
true maximal growth factor for a given pattern is bounded above by the smallest
recorded bound among all such valid upper bounds.

Because the recurrence is monotone with respect to the placement of delayed doubling,
no other admissible configuration can exceed these extremal realizations.
Accordingly, combining the sharpest bounds obtained across different initiating
rows yields the tightest global envelope on admissible local growth.

\[
\begin{array}{c|c|c}
\text{Pattern} & \text{Max Ratio} & \text{Associated coefficient count} \\
\hline
s,d & 3 & 2 \\
d,d & 123/33 & 1 \\
s,d,s & 341/131 & 1 \\
d,d,s & 329/123 & 1
\end{array}
\]

Consequently, the maximal geometric growth rate of a $5$--cycle satisfies
\begin{equation}
\Bigl(
3^2 \cdot (123/33) \cdot (341/131) \cdot (329/123)
\Bigr)^{1/5} 
\approx 2.976 < 3.
\label{eq:rho-5}
\end{equation}

% This bound directly feeds into Theorem~\ref{thm:maximal-growth} and ensures that
This bound directly feeds into Theorem ??? and ensures that
the per-symbol growth of any admissible infinite sequence under \OEIS{A186009}
is strictly dominated by $3$, as explained in the main text.

\paragraph{Conclusion.}
Every admissible growth pattern of \OEIS{A186009} is constrained by the grammar
of \OEIS{A022921}. Infinite concatenations of 5-cycles, though forbidden, saturate
local growth constraints and provide a strict upper envelope. Equation~\eqref{eq:rho-5}
confirms that even this extremal envelope has geometric growth rate $< 3$, ensuring
subdominance relative to the denominator $2^{B(i)}$.

% ============================================================
\newpage
\section{Grammar–Drift–Block Framework}
\label{sec:grammar-drift-block}
% ============================================================

\begin{table}[htbp]
\centering
\caption{Minimum block size, $M$, for logarithmic convexity of $a(n) = A186009(n)$}
\label{tab:A186009}
\vspace{4pt}
\begin{tabular}{@{}rccrr@{}}
\toprule
$M$ & $a_{n} a_{n+2M} \ge a_{n+M}$ & {S/D} & $n$ & $a(n)$ \\
\midrule
1 & $a_{ 1} a_{ 3} \ge a_{ 2} a_{ 2}$ & S &  1 &            1 \\
1 & $a_{ 2} a_{ 4} \ge a_{ 3} a_{ 3}$ & S &  2 &            1 \\
2 & $a_{ 3} a_{ 7} \ge a_{ 5} a_{ 5}$ & S &  3 &            1 \\
1 & $a_{ 4} a_{ 6} \ge a_{ 5} a_{ 5}$ & D &  4 &            2 \\
2 & $a_{ 5} a_{ 9} \ge a_{ 7} a_{ 7}$ & S &  5 &            3 \\
1 & $a_{ 6} a_{ 8} \ge a_{ 7} a_{ 7}$ & D &  6 &            7 \\
1 & $a_{ 7} a_{ 9} \ge a_{ 8} a_{ 8}$ & S &  7 &           12 \\
3 & $a_{ 8} a_{14} \ge a_{11} a_{11}$ & D &  8 &           30 \\
1 & $a_{ 9} a_{11} \ge a_{10} a_{10}$ & D &  9 &           85 \\
2 & $a_{10} a_{14} \ge a_{12} a_{12}$ & S & 10 &          173 \\
1 & $a_{11} a_{13} \ge a_{12} a_{12}$ & D & 11 &          476 \\
1 & $a_{12} a_{14} \ge a_{13} a_{13}$ & S & 12 &          961 \\
4 & $a_{13} a_{21} \ge a_{17} a_{17}$ & D & 13 &         2652 \\
1 & $a_{14} a_{16} \ge a_{15} a_{15}$ & D & 14 &         8045 \\
3 & $a_{15} a_{21} \ge a_{18} a_{18}$ & S & 15 &        17637 \\
1 & $a_{16} a_{18} \ge a_{17} a_{17}$ & D & 16 &        51033 \\
2 & $a_{17} a_{21} \ge a_{19} a_{19}$ & S & 17 &       108950 \\
1 & $a_{18} a_{20} \ge a_{19} a_{19}$ & D & 18 &       312455 \\
1 & $a_{19} a_{21} \ge a_{20} a_{20}$ & S & 19 &       663535 \\
4 & $a_{20} a_{28} \ge a_{24} a_{24}$ & D & 20 &      1900470 \\
1 & $a_{21} a_{23} \ge a_{22} a_{22}$ & D & 21 &      5936673 \\
2 & $a_{22} a_{26} \ge a_{24} a_{24}$ & S & 22 &     13472296 \\
1 & $a_{23} a_{25} \ge a_{24} a_{24}$ & D & 23 &     39993895 \\
1 & $a_{24} a_{26} \ge a_{25} a_{25}$ & S & 24 &     87986917 \\
4 & $a_{25} a_{33} \ge a_{29} a_{29}$ & D & 25 &    257978502 \\
1 & $a_{26} a_{28} \ge a_{27} a_{27}$ & D & 26 &    820236724 \\
3 & $a_{27} a_{33} \ge a_{31} a_{31}$ & S & 27 &   1899474678 \\
1 & $a_{28} a_{30} \ge a_{29} a_{29}$ & D & 28 &   5723030586 \\
2 & $a_{29} a_{33} \ge a_{31} a_{31}$ & S & 29 &  12809477536 \\
1 & $a_{30} a_{32} \ge a_{31} a_{31}$ & D & 30 &  38036848410 \\
1 & $a_{31} a_{33} \ge a_{32} a_{32}$ & S & 31 &  84141805077 \\
4 & $a_{32} a_{40} \ge a_{36} a_{36}$ & D & 32 & 248369601964 \\
1 & $a_{33} a_{35} \ge a_{34} a_{34}$ & D & 33 & 794919136728 \\
\bottomrule
\end{tabular}
\end{table}

This section formalises the interaction between the combinatorial
structure induced by $A020914(n)$ and the induced growth dynamics
of $C(n) = \sum_{k=0}^n N(k)$.

The central idea is that admissible division words generate a
finite-state grammar, which induces a bounded drift automaton.
This in turn yields a uniform block inequality governing residual decay.

% ============================================================
\clearpage
\bibliographystyle{unsrtnat}
\bibliography{src/latex/references}

\end{document}